% !TEX root = ../main.tex
\clearpage
\chapter{Appendix}
\label{ch:appendix}

% Demote headings one level under Chapter 8 (sections become 8.1.x).
\renewcommand{\dissertationSubheading}[2][]{%
  \phantomsection
  \subsection{#2}%
  \if\relax\detokenize{#1}\relax\else\label{#1}\fi
}
\renewcommand{\dissertationSubsubheading}[2][]{%
  \phantomsection
  \subsubsection{#2}%
  \if\relax\detokenize{#1}\relax\else\label{#1}\fi
}

\section{Evaluation Pipeline and Reproducibility}
\label{ch:appendix-eval}

This appendix documents the evaluation pipeline behind Chapter~\ref{ch:results}: its stages, the extraction design, the scaling tables, the account-type classification, the log of every timed run, the default parameters, the proxies and a glossary. Paths are given relative to the folder \texttt{1-Dissertation-Works/ERC20-Allowance-PageRank-Contract-Reputation/experiments} of the public repository \url{https://github.com/abitocodes/phd_works}.

\dissertationSubheading[sec:pipeline-stages]{Pipeline stages}

The evaluation runs in six stages. Stages 2 to 4 read the logs of the observation window, 1~October~2023 to 30~June~2026 (UTC, inclusive at both ends); stage 5 adds July to September~2026; stages 1 and 6 read no logs.

\begin{enumerate}
\item Plan and configuration: The analysis plan (\texttt{docs/analysis\_plan.md}), the configuration (\texttt{config/contract\_reputation.yaml}), the materialization query (\texttt{sql/01\_materialize\_logs.sql}) and the BigQuery runner (\texttt{scripts/run\_bq.py}) were committed together in \texttt{43d8045}, before any log of the study was extracted. The plan fixes the cohorts, the scores, the proxies and labels, the statistics and rule~A for the coupled operator.
\item Scan~A and the cohort: \texttt{scripts/run\_bq.py} dry-runs every query, checks the estimate against the budget in the configuration, runs the query and writes the billed bytes to \texttt{data/bq\_ledger.json}. Scan~A materializes the observation window (Appendix~\ref{sec:bq-extraction}). \texttt{sql/02\_spender\_candidates.sql} lists the spenders that received a non-zero approval from at least three distinct owners in the window, \texttt{scripts/check\_account\_code.py} reads the account type of each (Appendix~\ref{sec:account-type-classification}), and \texttt{scripts/build\_cohort.py} keeps the contracts as the matched cohort and loads their addresses into BigQuery.
\item Graph tables, proxies and labels: \texttt{sql/03\_ego\_events.sql} extracts the ego network of the cohort, and \texttt{sql/04\_anchor\_pairs.sql} aggregates it into the pair tables of both graphs at $T_{\text{obs}}$ and at $t_1$. \texttt{sql/05\_proxies\_same\_window.sql} computes the transfer and Sybil-stability proxies and \texttt{sql/06\_holdout\_labels.sql} the labels of W0. \texttt{sql/07\_gmx\_accounts.sql} and \texttt{sql/09\_gmx\_contract\_closes.sql} select the GMX~V2 closes of contract accounts, which \texttt{scripts/decode\_gmx\_events.py} decodes with the event ABI kept in \texttt{data/0-gmx-abi-for-decoding-logs/}. \texttt{sql/08\_export\_ids.sql} replaces addresses by integer node ids, \texttt{scripts/download\_graph\_tables.sh} downloads the tables in parquet parts, and \texttt{sql/10\_describe.sql} counts the events and allowance amounts described in Appendix~\ref{ch:appendix-data}.
\item Same window and W0: \texttt{scripts/run\_same\_window.py} scores every method at $T_{\text{obs}}$ and computes the same-window checks, the contrasts and the secondary part of rule~A. \texttt{scripts/run\_holdout.py} scores every method at $t_1$ and evaluates the spender and trader cohorts of W0 together with the primary part of rule~A.
\item Registration and scan~B: The registration of W1 (\texttt{docs/registration\_w1.md}) was committed in \texttt{0d3a057} with the configuration status set to \texttt{registered}. \texttt{scripts/extract\_w1.py} refuses to run unless both files are committed without local changes and the status is \texttt{registered}; it writes their SHA-256 hashes and the commit to \texttt{registered-w1/extraction\_manifest.json} and then runs scan~B, the ego extraction of July--September and the GMX account query of that window. The W1 labels and closes follow from \texttt{sql/06} and \texttt{sql/09}, and \texttt{scripts/run\_registered.py} evaluates rule~B, its two sensitivity analyses and the comparison on outcomes built from neither edge type, in the order of the registration.
\item After the labels: \texttt{scripts/run\_supplementary.py} computes the spender contrasts made after the labels were known, \texttt{scripts/run\_robustness.py} the top-token subgraphs (with \texttt{sql/11\_top\_tokens.sql}) and the sample-definition stages, \texttt{scripts/sybil\_model.py} the Sybil cost model, and \texttt{scripts/run\_benchmark.py} the timed runs. \texttt{scripts/export\_latex.py} writes the tables of \texttt{results/tables/} and the figures of \texttt{Figures/generated/} from the summaries, and \texttt{scripts/publish\_results.py} copies the summaries to \texttt{data/3-published-results-for-thesis/} with a digest that names the source of every number and a list of checksums.
\end{enumerate}

The order of these steps is recorded in the commit history, in the cost ledger and in the extraction manifest of W1. The plan was committed at 13:41~UTC on 9~October~2026; the ledger, which records each job when it completes, lists the first job of scan~A at 13:43~UTC. The other queries and scripts of stages 2 to 4 entered the repository with the registration in \texttt{0d3a057}, after they had run, together with the summaries of the same window and of W0 and with the scripts of stage 5, which had not yet run. That commit was made at 16:49:32~UTC on 9~October~2026, which the commit records in Korea Standard Time as 01:49:32 on 10~October. \texttt{scripts/extract\_w1.py} wrote its manifest one second later, before its first query, and the manifest records that no override was used. The digest \texttt{results\_digest.md} records the code revision at which the summaries were published, \texttt{8d8b568}, and names the scripts that differed from that revision when the digest was written. Commit \texttt{8d8b568} added or extended the scripts of stage 6 and changed \texttt{scripts/run\_holdout.py} only in how it groups the contrasts of W0, which the summary of W0 records as leaving every value unchanged. The benchmark ran with \texttt{scripts/run\_benchmark.py} as committed in \texttt{8d8b568}, the revision the digest records; the script has not changed since. The timed runs were made on an otherwise idle machine with an Intel Core Ultra~9 185H processor (16 cores, 22 logical processors) and 64~GB of memory, under Windows~11 (build 26200), with Python 3.11.9, NumPy 2.3.2 and SciPy 1.16.1. \texttt{benchmark/benchmark.json} records the platform (\texttt{Windows-10-10.0.26200-SP0}, the string that Python reports for this build), the processor (Intel64 Family~6 Model~170 Stepping~4) and the Python version.

The public repository \url{https://github.com/abitocodes/phd_works} contains the evaluation source in \texttt{1-Dissertation-Works/ERC20-Allowance-PageRank-Contract-Reputation/experiments}: the queries (\texttt{sql/}), the scripts (\texttt{scripts/}), the tests (\texttt{tests/}), the versioned configuration (\texttt{config/contract\_reputation.yaml}), the plan and the registration (\texttt{docs/}), the extracted tables, the summaries and the cost ledger. Every data file is split into parquet parts below 90\,MB so that it can be committed, and a JSON file beside each table records the BigQuery table it came from, its row count and the time of the download. The summaries behind the tables are in \texttt{data/3-published-results-for-thesis/}, together with copies of the configuration, the plan and the registration, the digest \texttt{results\_digest.md}, which gives the summary file and key of every number, and \texttt{SHA256SUMS.txt}. Any number in a table can thus be checked against its summary without re-running the pipeline.

The scoring and statistics code is checked offline: \texttt{tests/test\_scoring.py} and \texttt{tests/test\_stats.py} compare the solver, the edge weights and the paired bootstrap on small synthetic graphs with a reference implementation of the same definitions, with no cloud access. Real extraction needs Google Cloud credentials and a billing project, and its scans read only the public Arbitrum One log table.

\dissertationSubheading[sec:bq-extraction]{BigQuery extraction design}

The public table \texttt{bigquery-public-data.goog\_blockchain\_arbitrum\_one\_us.logs} is read once for each window. Scan~A runs \texttt{sql/01\_materialize\_logs.sql} over 1~October~2023 to 30~June~2026 and keeps four kinds of log in one project table, \texttt{contract\_rep.logs\_obs}, partitioned by month: ERC-20 \texttt{Approval} and \texttt{Transfer} logs with exactly three topics, identified by their standard topic-0 hashes \parencite{vogelsteller2015erc20}; GMX~V2 \texttt{PositionDecrease} logs of the protocol's event emitter; and the \texttt{Borrow}, \texttt{Repay} and \texttt{LiquidationCall} logs of the Aave~V3 pool, which no analysis of this thesis uses. Scan~B runs the same query over July to September~2026 into \texttt{contract\_rep.logs\_w1}. Every later query reads these project tables, so the cost of the chain-wide scan is paid once per window. Appendix~\ref{ch:appendix-sql} shows the three queries on which the graphs rest.

The ego extraction (\texttt{sql/03\_ego\_events.sql}) keeps an approval when a contract of the matched cohort is its owner or its spender, and a transfer when a cohort contract is its sender or its recipient. Each graph is therefore the ego network of the cohort, and the size of the extract follows the cohort and not the chain. Amounts are raw base units, the \texttt{uint256} value read as a double. Zero-value transfers are dropped, while zero-value approvals are kept because they delete an allowance. For the observation window the two ego tables hold $389{,}866{,}882$ approval logs and $1{,}538{,}404{,}067$ transfer logs.

The pair tables (\texttt{sql/04\_anchor\_pairs.sql}) are the edge lists of the two graphs at one anchor. They are built at $T_{\text{obs}}$ and at $t_1$ from the events up to the anchor, with the time decay anchored there, and each pair carries the weights of every method (Appendix~\ref{sec:sql-pairs}). At $T_{\text{obs}}$ they hold $22{,}321{,}470$ allowance pairs and $59{,}737{,}617$ transfer pairs, and at $t_1$ $21{,}550{,}791$ and $58{,}065{,}784$. The freeze of W1 coincides with $T_{\text{obs}}$, so W1 reuses the pair tables of $T_{\text{obs}}$. \texttt{sql/08\_export\_ids.sql} numbers the $25{,}305{,}099$ addresses that occur in a pair table of either anchor or in the cohort. Only the id-coded pair tables, the node table, the proxies and the labels are downloaded; the event tables stay in BigQuery.

Every job is listed in \texttt{data/bq\_ledger.json} with its query, parameters, estimated and billed bytes and cost. The ledger holds 17 jobs that billed $5{,}226{,}856{,}710{,}144$ bytes ($4.75$\,TiB), or USD~29.71 at the on-demand rate of USD~6.25 per TiB recorded in the ledger, against a budget of USD~50 fixed in the configuration. Scan~A alone billed $3.29$\,TiB (USD~20.55) and scan~B $0.26$\,TiB (USD~1.64). When the registration was committed, the ledger stood at USD~26.20.

\dissertationSubheading[sec:scaling-appendix]{Scaling tables}

Table~\ref{tab:benchmark-scaling} in Section~\ref{sub:runtime-scaling} reports the runtime of EndorseRank and AWP on deterministic SHA256 subsamples of the matched cohort, measured under the timing protocol of Section~\ref{sec:computational-benchmark}; this section describes how its stages were drawn. The contracts are ordered by the SHA256 hash of the seed string \texttt{contract-benchmark-v1} joined to the address, and the first $n$ of them form a stage, for $n=1{,}000$, $2{,}000$, $5{,}000$, $10{,}000$ and $20{,}000$; the last row is the full cohort, whose measurement is the one in Table~\ref{tab:benchmark-runtime}, reused. A stage keeps every extracted edge with an end point among its contracts, so $|V|$ and $|E|$ of both graphs grow with $n$, from $194{,}084$ allowance and $1{,}385{,}509$ transfer edges at $n=1{,}000$ to the full graphs with $22{,}321{,}470$ and $59{,}737{,}617$, and the table gives both counts for every stage. The stages are the evaluation sets of Table~\ref{tab:robustness-sample-size}, so the edge counts in the two tables agree.

The ratio of transfer edges to allowance edges is not constant across the stages. It is highest at the smallest stage and settles near its full-cohort value from $n=5{,}000$ on; the table gives it for every stage, and Section~\ref{sub:runtime-scaling} discusses it. The table therefore shows whether the ratio of the two runtimes follows the ratio of the edge counts from stage to stage. No stage is larger than the cohort graph, so the table says nothing about graphs beyond it.

\dissertationSubheading[sec:account-type-classification]{Account-type classification}

Only contracts are scored and evaluated, so every address that could enter a cohort was classified by its code first. \texttt{scripts/check\_account\_code.py} calls \texttt{eth\_getCode} for each address at one block, fetched once at the start of a run and written to the output with the time of the read. An address with non-empty code is a contract unless the code is an EIP-7702 delegation designator, the bytes \texttt{0xef0100} followed by a 20-byte address \parencite{buterin2024eip7702}; such an address, like one without code, is an externally owned account (EOA). Table~\ref{tab:account-types} lists the four runs. The candidates of the matched cohort were read from the public endpoint \url{https://arb1.arbitrum.io/rpc} until it limited the request rate after $20{,}000$ addresses; the other $9{,}022$ were read at the same block from \url{https://arbitrum-one.public.blastapi.io}, which also served the three later runs.

\begin{table}[htbp]
\centering
\caption[Account types read with \texttt{eth\_getCode}.]{Account types read with \texttt{eth\_getCode} on Arbitrum One, one block per run. Contract: non-empty code that is not a delegation designator. No code and EIP-7702 (a delegation designator) are both externally owned accounts. The W0 run read the spenders that were not cohort candidates, and the W1 run the addresses that no earlier run had classified. Sources: the JSON files beside the parquet outputs in \texttt{data/1-raw-blockchain-extracts/cohort/} and \texttt{gmx/}.}
\label{tab:account-types}
\fitwidth{%
\begin{tabular}{lrrrrrr}
\toprule
Address set & Block & Read at (UTC) & Addresses & Contract & No code & EIP-7702 \\
\midrule
Cohort candidates: non-zero approvals from $\geq 3$ owners, observation window & 513{,}216{,}506 & 9 Oct 2026, 13:55 & 29{,}022 & 27{,}844 & 1{,}121 & 57 \\
GMX~V2 accounts with $\geq 3$ closes in the observation window & 513{,}219{,}999 & 9 Oct 2026, 14:10 & 40{,}307 & 7{,}633 & 28{,}075 & 4{,}599 \\
W0 spenders with a positive allowance at $t_1$, not cohort candidates & 513{,}233{,}527 & 9 Oct 2026, 15:08 & 8{,}224 & 7{,}792 & 413 & 19 \\
W1 spenders and GMX~V2 accounts with $\geq 3$ closes in W1, not read before & 513{,}257{,}476 & 9 Oct 2026, 16:52 & 1{,}153 & 350 & 568 & 235 \\
\bottomrule
\end{tabular}
}
\end{table}

The matched cohort consists of the contracts among the candidates. The W0 spender cohort draws on the first and third runs: of the spenders with a positive latest allowance at $t_1$, $33{,}101$ are contracts, $431$ have no code and $24$ carry a delegation designator. The W1 spender cohort and the two trader cohorts take their account types from all four runs. Every read refers to the chain on 9~October~2026, so an address is classified by the code it held on that day.

\dissertationSubheading[sec:benchmark-run-logs]{Benchmark run logs}

Table~\ref{tab:benchmark-runs} lists the wall-clock seconds of every timed solver call of the benchmark, in the order the runs were made, with their mean and standard deviation: the four scores on the full graphs at $d=0.85$, EndorseRank and AWP at $d=0.75$ and $d=0.95$, and the two scores at every scaling stage below the full cohort. Each configuration had one untimed warm-up before its five timed runs, and peak memory was traced on the warm-up alone. Each graph and setting was timed once, and that measurement is reused wherever the same graph and setting appear: the $d=0.85$ rows of Table~\ref{tab:robustness-damping} and the full-cohort row of Table~\ref{tab:benchmark-scaling} repeat the first block of Table~\ref{tab:benchmark-runs}. Across all configurations the standard deviation of the five runs is between 1.5\% and 5.1\% of the mean.

% Generated by experiments/scripts/export_latex.py -- do not edit by hand.
% Source: experiments/data/2-processed-tables-and-evaluations/revised-usd/benchmark/benchmark.json
\begin{table}[htbp]
\centering
\caption[Every timed solver call of the benchmark.]{Wall-clock seconds of every timed solver call of the benchmark (matched cohort, $n=15{,}107$ contracts), in the order the runs were made, with their mean and standard deviation. Each configuration had one untimed warm-up before its timed runs, and each graph and setting was timed once: the $d=0.85$ rows of Table~\ref{tab:robustness-damping} and the full-cohort row of Table~\ref{tab:benchmark-scaling} are the first block, reused. The stages are the subsamples of Table~\ref{tab:benchmark-scaling}.}
\label{tab:benchmark-runs}
\fitwidth{%
\begin{tabular}{lrrrrrrr}
\toprule
Method & Run 1 & Run 2 & Run 3 & Run 4 & Run 5 & Mean & SD \\
\midrule
\multicolumn{8}{l}{\emph{Full graphs, $d=0.85$}} \\
EndorseRank & 9.570 & 11.063 & 10.497 & 9.811 & 10.530 & 10.294 & 0.601 \\
AWP & 21.139 & 20.446 & 24.122 & 24.492 & 21.637 & 22.367 & 1.825 \\
C-PR ($\lambda=0.5$) & 33.813 & 39.576 & 35.439 & 35.793 & 36.187 & 36.162 & 2.112 \\
S-PR & 38.769 & 36.970 & 29.310 & 35.214 & 31.848 & 34.422 & 3.833 \\
\midrule
\multicolumn{8}{l}{\emph{Full graphs, $d=0.75$}} \\
EndorseRank & 6.045 & 5.963 & 6.003 & 5.984 & 5.927 & 5.984 & 0.044 \\
AWP & 13.208 & 16.968 & 16.667 & 15.870 & 13.469 & 15.236 & 1.781 \\
\midrule
\multicolumn{8}{l}{\emph{Full graphs, $d=0.95$}} \\
EndorseRank & 31.701 & 32.171 & 30.886 & 30.574 & 30.836 & 31.233 & 0.673 \\
AWP & 58.498 & 56.896 & 59.713 & 64.490 & 65.212 & 60.962 & 3.697 \\
\midrule
\multicolumn{8}{l}{\emph{Stage $n=1{,}000$}} \\
EndorseRank & 0.425 & 0.489 & 0.585 & 0.683 & 0.719 & 0.580 & 0.125 \\
AWP & 1.892 & 1.408 & 1.870 & 1.758 & 1.458 & 1.677 & 0.229 \\
\midrule
\multicolumn{8}{l}{\emph{Stage $n=2{,}000$}} \\
EndorseRank & 0.944 & 0.956 & 0.925 & 0.932 & 0.937 & 0.939 & 0.012 \\
AWP & 3.748 & 4.876 & 4.133 & 4.811 & 3.792 & 4.272 & 0.543 \\
\midrule
\multicolumn{8}{l}{\emph{Stage $n=5{,}000$}} \\
EndorseRank & 2.293 & 2.420 & 2.489 & 2.555 & 2.723 & 2.496 & 0.159 \\
AWP & 8.713 & 9.861 & 9.895 & 8.051 & 8.453 & 8.995 & 0.840 \\
\midrule
\multicolumn{8}{l}{\emph{Stage $n=10{,}000$}} \\
EndorseRank & 6.253 & 6.375 & 6.065 & 6.285 & 6.119 & 6.220 & 0.126 \\
AWP & 17.640 & 17.362 & 17.160 & 16.124 & 15.816 & 16.821 & 0.802 \\
\bottomrule
\end{tabular}
}
\end{table}


\dissertationSubheading[sec:parameter-summary]{Parameter summary}

Table~\ref{tab:parameter-summary} collects the default parameters of the evaluation. Most of them are keys of the configuration (Appendix~\ref{sec:config-keys}) and were set with the plan before any log was extracted. The registration of W1 added the margin of rule~B and the $2{,}000$ resamples of the trader comparisons, and the block of the code check was fixed when the check ran.

\begin{table}[htbp]
\centering
\caption{Default parameters used in the evaluation.}
\label{tab:parameter-summary}
\fitwidth{%
\begin{tabular}{lll}
\toprule
Parameter & Value & Role \\
\midrule
PageRank damping $d$ & $0.85$; sweep $0.75$, $0.95$ & Teleport mixture \\
Convergence tolerance $\varepsilon$ & $10^{-8}$ ($L_1$) & Power-iteration stop \\
Max iterations $I_{\max}$ & $300$ & Power-iteration cap \\
Logistic $k$ & $0.01$ & Decay slope (AWP, EndorseRank, transfer layer of C-PR) \\
Logistic $t_0$ & $180$ days & Decay midpoint (AWP, EndorseRank, transfer layer of C-PR) \\
Value transform $b$ & $1$ (raw base units) & Value weight $V(z)$ of AWP and EndorseRank \\
Restart distribution & $\propto X_u$; uniform & AWP; EndorseRank, C-PR \\
Coupling weight $\lambda$ & $0.5$ (primary); $0.25$, $0.75$ & C-PR allowance-layer weight, fixed in the plan \\
S-PR restart vector & C-PR ($\lambda=1$) on the transfer-layer nodes, renormalized & Endorsement-seeded teleportation \\
Cohort rule & non-zero approvals from $\geq 3$ distinct owners; contract code & Matched cohort \\
Account type & \texttt{eth\_getCode} at block 513{,}216{,}506 (cohort) & Contract or EOA \\
Observation window & 2023-10-01 -- 2026-06-30 & UTC inclusive; $T_{\text{obs}}$ anchors the decay \\
W0 freeze / labels & 2026-03-31 / 2026-04-01 -- 2026-06-30 & Spender and trader cohorts; rule~A \\
W1 freeze / labels & 2026-06-30 / 2026-07-01 -- 2026-09-30 & Registered window; rule~B, margin $\delta=0.01$ \\
Trader rule & contract code, node at the freeze, $\geq 3$ closes in the label window & GMX~V2 trader cohorts \\
Bootstrap resamples / seed & $400$ / $42$; $2{,}000$ / $42$ & Paired contract bootstrap; trader labels of Table~\ref{tab:neutral-label-grid} \\
Timed runs / warm-up & $5$ / $1$ & Runtime protocol \\
Scaling seed & \texttt{contract-benchmark-v1} & Deterministic SHA256 subsamples \\
Top-token lists & $20$ by summed raw amount; $20$ by rows & Token subgraphs (Table~\ref{tab:robustness-tokens}) \\
BigQuery budget & USD~50 & Checked before every query \\
\bottomrule
\end{tabular}
}
\end{table}

\dissertationSubheading[sec:proxy-appendix]{Proxy operational definitions (reference)}

The proxies are defined in Section~\ref{sec:proxy-formulas}; the list below names each one with its equation and states in one line what it counts for a cohort contract $v$.

\begin{itemize}
\item Transfer family: the in-degree (Equation~\eqref{eq:in-degree}) counts the distinct senders of transfers to $v$, self-transfers included; the in-value (Equation~\eqref{eq:in-value}) sums their raw amounts across tokens.
\item Allowance family: the in-approve degree (Equation~\eqref{eq:in-approve-degree}) counts the owners with a positive latest allowance to the spender $v$; the in-approve value (Equation~\eqref{eq:in-approve-value}) sums those allowances across owners and tokens.
\item Sybil-adjusted stability family: the inbound counterparty ratio (Equation~\eqref{eq:counterparty-ratio}) divides the distinct other senders by the transfers received from them; the transfer tenure (Equation~\eqref{eq:tenure}) is the span in days of the transfers attributed to $v$; the active months (Equation~\eqref{eq:active-months}) count the calendar months with such transfers.
\end{itemize}

\dissertationSubheading[sec:glossary]{Glossary}

The terms below are used throughout the thesis with the meanings given here.

\begin{description}
\item[Account type] Contract or externally owned account, read with \texttt{eth\_getCode} at one block (Appendix~\ref{sec:account-type-classification}).
\item[Activeness] AWP's restart weight of an address: over the addresses it has sent to, the sum of the time weights of its latest transfer to each.
\item[Allowance] An ERC-20 authorization that lets a spender move tokens on an owner's behalf.
\item[Alignment] Rank agreement between a score and a proxy, measured by Spearman's $\rho$ on mid-ranks and by Kendall's $\tau$; every reported $\tau$ is the tie-corrected $\tau_b$.
\item[Construct overlap] A ranking method and a proxy sharing a data source, which raises their agreement; such agreement is a check of the intended construct and not evidence from outside it.
\item[C-PR end points] C-PR at $\lambda=1$ and at $\lambda=0$, which walk the raw-amount allowance layer or transfer layer alone with uniform restarts.
\item[Ego network] The events with a cohort contract at one end: the owner or spender of an approval, the sender or recipient of a transfer.
\item[Endorsement graph] Directed graph of positive latest allowances from owner to spender, the input of EndorseRank and the allowance layer of C-PR.
\item[Externally owned account (EOA)] An address without code or with an EIP-7702 delegation designator. EOAs appear in the graphs as owners, senders and recipients and are never scored or evaluated.
\item[Matched cohort] The primary sample: the contracts that received a non-zero \texttt{Approval} from at least three distinct owners in the observation window, fixed from the approval logs and the code check before any transfer was aggregated (Section~\ref{sub:cohorts}). Membership does not depend on either graph at $T_{\text{obs}}$, so a cohort contract can lie outside one graph or both.
\item[Paired bootstrap] Resampling of contracts with replacement through a single index matrix shared by every coefficient, so that differences between coefficients can be tested on the same resamples.
\item[PageRank] Spectral ranking of nodes by the stationary distribution of a damped random walk.
\item[Proxy] An observable statistic of a contract against which a ranking is checked for the construct it is built to measure.
\item[Registered window (W1)] Scores frozen at 30~June~2026 and labels counted from July to September~2026, with rule~B and the comparison on outcomes built from neither edge type registered before the logs of those months were extracted (Section~\ref{sub:fresh-holdout}).
\item[Scaling stage] A deterministic SHA256 subsample of the matched cohort together with the extracted edges incident to its contracts.
\item[Spender cohort] Every contract spender with at least one positive latest allowance at the freeze in the extracted approval data: at $t_1$ (31~March~2026, 23:59:59 UTC) for W0 and at 30~June~2026, 23:59:59 UTC, for W1.
\item[Teleport] The jump that PageRank makes with probability $1-d$, to a node drawn from the restart distribution: uniform for EndorseRank and C-PR, weighted by sending activity for AWP.
\item[Trader cohort] The GMX~V2 accounts that hold contract code, are nodes of the graph at the freeze and close at least three positions in the label window.
\end{description}

% !TEX root = ../main.tex
\section{SQL Templates, Configuration, and Extended Notes}
\label{ch:appendix-sql}

To aid reproduction, this appendix shows the three SQL patterns on which every graph of the thesis rests, the configuration keys, and notes on complexity and on the conventions of the result tables. The listings are abridged; the full files are in \texttt{sql/} of the experiments folder (Appendix~\ref{ch:appendix-eval}). Placeholders in double underscores, such as \texttt{\_\_DEST\_\_}, are replaced by \texttt{scripts/run\_bq.py} before a query runs, and names that start with \texttt{@} are query parameters.

\dissertationSubheading[sec:sql-materialize]{Log materialization pattern}

Scan~A and scan~B run the same query, \texttt{sql/01\_materialize\_logs.sql}, over different windows. It reads the public log table once, tags each log by kind and keeps only the four kinds the study needs, in a project table partitioned by month and clustered by kind. With the GMX and Aave branches shortened, it reads:

\begin{verbatim}
CREATE TABLE `__DEST__`
PARTITION BY TIMESTAMP_TRUNC(block_timestamp, MONTH)
CLUSTER BY kind
AS
WITH src AS (
  SELECT block_timestamp, block_number, log_index,
         LOWER(address) AS address, topics, data
  FROM `bigquery-public-data.goog_blockchain_arbitrum_one_us.logs`
  WHERE block_timestamp >= TIMESTAMP(@start_ts)
    AND block_timestamp < TIMESTAMP(@end_ts)
),
tagged AS (
  SELECT *,
    CASE
      WHEN ARRAY_LENGTH(topics) = 3 AND topics[OFFSET(0)] =
        '0x8c5be1e5ebec7d5bd14f71427d1e84f3dd0314c0f7b2291e5b200ac8c7c3b925'
        THEN 'A'                                  -- ERC-20 Approval
      WHEN ARRAY_LENGTH(topics) = 3 AND topics[OFFSET(0)] =
        '0xddf252ad1be2c89b69c2b068fc378daa952ba7f163c4a11628f55a4df523b3ef'
        THEN 'T'                                  -- ERC-20 Transfer
      WHEN address = '0xc8ee91a54287db53897056e12d9819156d3822fb'
           AND ... THEN 'G'                       -- GMX V2 PositionDecrease
      WHEN address = '0x794a61358d6845594f94dc1db02a252b5b4814ad'
           AND ... THEN 'V'                       -- Aave V3 pool events
    END AS kind
  FROM src
)
SELECT
  block_timestamp, block_number, log_index, kind,
  FROM_HEX(SUBSTR(address, 3)) AS emitter,
  IF(kind IN ('A', 'T'), FROM_HEX(SUBSTR(topics[OFFSET(1)], 27)), NULL) AS a,
  IF(kind IN ('A', 'T'), FROM_HEX(SUBSTR(topics[OFFSET(2)], 27)), NULL) AS b,
  IF(kind IN ('A', 'T'), SAFE.FROM_HEX(SUBSTR(data, 3)), NULL) AS val,
  IF(kind IN ('G', 'V'), topics, NULL) AS raw_topics,
  IF(kind IN ('G', 'V'), data, NULL) AS raw_data
FROM tagged
WHERE kind IS NOT NULL
\end{verbatim}

For an approval, \texttt{a} is the owner and \texttt{b} the spender; for a transfer, \texttt{a} is the sender and \texttt{b} the recipient. In both, \texttt{emitter} is the token contract and \texttt{val} the 32-byte amount. The two topic-0 hashes are the keccak256 hashes of \texttt{Approval(address,address,uint256)} and \texttt{Transfer(address,address,uint256)} \parencite{vogelsteller2015erc20}. Requiring exactly three topics keeps the ERC-20 logs and leaves out the ERC-721 logs, which share both hashes but index the token id as a fourth topic. Addresses are stored as 20-byte values, not as hexadecimal strings.

\dissertationSubheading[sec:sql-ego]{Ego-network extraction pattern}

\texttt{sql/03\_ego\_events.sql} reads the materialized table of one window and keeps the approvals and transfers with a cohort contract at one end. The table \texttt{contract\_rep.cohort} holds the 20-byte addresses of the matched cohort, loaded by \texttt{scripts/build\_cohort.py}. Abridged:

\begin{verbatim}
CREATE TEMP FUNCTION u256(v BYTES) AS (
  IFNULL((SELECT SUM(c * POW(256.0, LENGTH(v) - 1 - o))
          FROM UNNEST(TO_CODE_POINTS(v)) AS c WITH OFFSET o), 0.0)
);

CREATE OR REPLACE TABLE `dissertation-bq.contract_rep.approvals___SUFFIX__`
PARTITION BY TIMESTAMP_TRUNC(block_timestamp, MONTH)
AS
SELECT
  l.block_timestamp, l.block_number, l.log_index,
  l.emitter AS token_address,
  l.a AS owner,
  l.b AS spender,
  u256(l.val) AS value
FROM `dissertation-bq.contract_rep.__SRC__` AS l
WHERE l.kind = 'A'
  AND l.block_timestamp >= TIMESTAMP(@start_ts)
  AND l.block_timestamp < TIMESTAMP(@end_ts)
  AND (l.a IN (SELECT address FROM `dissertation-bq.contract_rep.cohort`)
       OR l.b IN (SELECT address FROM `dissertation-bq.contract_rep.cohort`));

-- transfers___SUFFIX__: the same with kind = 'T', from_address = a and
-- to_address = b, and zero amounts dropped:
--   AND LTRIM(TO_HEX(l.val), '0') != ''
\end{verbatim}

The function \texttt{u256} reads the 32-byte amount as a double in raw base units; no token decimals are applied anywhere in the pipeline. Approvals of zero are kept, because an approval of zero deletes an allowance, while zero-value transfers are dropped. The suffix is \texttt{ego} for the observation window and \texttt{w1ego} for July--September~2026.

\dissertationSubheading[sec:sql-pairs]{Latest-allowance and pair-table pattern}

\texttt{sql/04\_anchor\_pairs.sql} turns the ego events up to one anchor into the edge lists of both graphs. Abridged:

\begin{verbatim}
CREATE TEMP FUNCTION sig(age_days FLOAT64) AS
  (1.0 / (1.0 + EXP(0.01 * (age_days - 180.0))));
CREATE TEMP FUNCTION vt(x FLOAT64) AS
  (IF(x >= 40.0, 1.0, 2.0 / (1.0 + EXP(-x)) - 1.0));

CREATE OR REPLACE TABLE `dissertation-bq.contract_rep.allow_triples___TAG__` AS
SELECT token_address, owner, spender, value, block_timestamp,
       TIMESTAMP_DIFF(TIMESTAMP(@anchor_ts), block_timestamp, MICROSECOND)
         / 8.64e10 AS age_days
FROM (
  SELECT *, ROW_NUMBER() OVER (
           PARTITION BY token_address, owner, spender
           ORDER BY block_number DESC, log_index DESC) AS rn
  FROM `dissertation-bq.contract_rep.approvals_ego`
  WHERE block_timestamp >= TIMESTAMP(@start_ts)
    AND block_timestamp <= TIMESTAMP(@anchor_ts)
)
WHERE rn = 1 AND value > 0;

CREATE OR REPLACE TABLE `dissertation-bq.contract_rep.allow_pairs___TAG__` AS
SELECT owner, spender,
       COUNT(*) AS n_tokens,
       SUM(sig(age_days) * vt(value)) AS w_er,   -- EndorseRank
       SUM(value) AS w_raw,                      -- C-PR allowance layer
       MAX(sig(age_days)) AS max_sig,
       MAX(block_timestamp) AS last_ts
FROM `dissertation-bq.contract_rep.allow_triples___TAG__`
GROUP BY owner, spender;

-- transfer_pairs___TAG__: per (from_address, to_address) over transfers_ego
-- up to the anchor, with
--   SUM(sig(age_days) * vt(value)) AS w_awp    -- AWP
--   SUM(value * sig(age_days))     AS w_raw    -- C-PR transfer layer
--   MAX(sig(age_days))             AS max_sig  -- AWP activeness
\end{verbatim}

The function \texttt{sig} is the logistic decay of Equation~\eqref{eq:logistic-decay} with $k=0.01$ and $t_0=180$ days, and \texttt{vt} the value transform of Equation~\eqref{eq:value-transform} with $b=1$, set to one from $z=40$ on, where it equals one in double precision. Among the approvals of one token, owner and spender up to the anchor, the latest is the one with the largest block number and, within the block, the largest log index. It is kept only when its value is positive, so a revocation removes that token from the pair. Each pair row carries the weights of every method, so all scores at one anchor come from the same rows. The tag is \texttt{tobs} for $T_{\text{obs}}$ and \texttt{t1} for the freeze of W0.

\dissertationSubheading[sec:config-keys]{Configuration keys}

All of the pipeline's settings are kept in one versioned file, \texttt{config/contract\_reputation.yaml}, which the scripts read at start-up. Its groups of keys are:

\begin{itemize}
\item \texttt{chain\_id}, \texttt{chain\_name} and \texttt{bigquery}: chain 42161 (Arbitrum One); the project, dataset, location and public log table; the budget in USD; and the windows of scan~A and scan~B;
\item \texttt{period}, \texttt{holdout} and \texttt{registered\_window}: the observation window with its anchor $T_{\text{obs}}$, the freeze and label window of W0, and the freeze, label window, status and margin $\delta$ of W1;
\item \texttt{events}, \texttt{gmx\_arbitrum} and \texttt{aave\_arbitrum}: the topic-0 hashes of \texttt{Approval} and \texttt{Transfer}; the GMX~V2 event emitter, the \texttt{EventLog1} topic, the \texttt{PositionDecrease} hash, the liquidation handler, the order type of a liquidation and the minimum of three closes; and the Aave~V3 pool with its event topics;
\item \texttt{accounts} and \texttt{cohort}: the RPC endpoint and the delegation prefix \texttt{0xef0100} of the code check, and the minimum of three distinct owners;
\item \texttt{reputation}: damping, tolerance, iteration cap, decay $k$ and $t_0$, value parameter $b$, the primary and sensitivity values of $\lambda$, and the two parts of rule~A;
\item \texttt{alignment}, \texttt{benchmark} and \texttt{robustness}: the bootstrap resamples and seed; the warm-up and timed runs and the scaling seed; the damping values and the number of top tokens;
\item \texttt{storage}: the largest size of a data file part, 90\,MB.
\end{itemize}

Changing the damping factor, the decay or the value parameter requires running the evaluation again, the damping, top-token and sample-definition checks included, and exporting the tables and figures again from the same summaries, so that every reported number matches its machine-readable source. \texttt{scripts/extract\_w1.py} records the SHA-256 hash of this file at the start of scan~B (Appendix~\ref{sec:map-artifacts}).

\dissertationSubheading[sec:complexity-notes]{Complexity notes}

Write $m=|E|$. In BigQuery the latest allowances come from a window function over the $R_a$ approval rows, which sorts within each token, owner and spender, in $O(R_a\log R_a)$ time at most; the transfer pairs are one aggregation over the $R_t$ transfer rows, in $O(R_t)$. Locally, PageRank over $I$ iterations costs $O(m+mI)$, the first term for building the sparse operator and the second for one sparse product per iteration (Appendix~\ref{sec:pr-sparsity}). C-PR walks the union of both layers, so one of its iterations costs as much as one on each layer, and S-PR solves the allowance layer first and then the transfer layer. The edge lists and score vectors take most of the memory, so the peak memory of a solve grows with $m$; in Table~\ref{tab:benchmark-runtime} the peak-memory ratio of AWP to EndorseRank is $2.33$ ($2{,}995.9$ against $1{,}285.0$~MB), at an edge ratio of $2.68$.

\dissertationSubheading[sec:interpretation-checklist]{Conventions for the result tables}

The tables of Chapter~\ref{ch:results} and of this appendix follow these conventions:

\begin{enumerate}
\item Every $\tau$ is Kendall's $\tau_b$ on raw scores and proxies, reported with its 95\% paired-bootstrap interval: $400$ resamples over contracts with seed 42, and $2{,}000$ for the trader comparisons on outcomes built from neither edge type, whose decision contrasts P and Q also carry 97.5\% intervals. Differences between coefficients come from the paired contrast tables, not from subtracting point estimates.
\item A cohort contract missing from a method's graph scores zero under that method. Keeping it as an isolated node instead is the sensitivity analysis of Table~\ref{tab:tie-sensitivity}.
\item Each contrast table says which of its contrasts were fixed in advance, by the plan (rule~A) or by the registration of W1 (rule~B and the comparison on outcomes built from neither edge type), and which were computed after the labels were known.
\item A runtime is the mean wall-clock time of five timed solver calls after one untimed warm-up. A call turns edge lists built in advance into scores: node indexing, sparse-matrix assembly, row normalization, restart vector and power iteration, without extraction or preprocessing. Every runtime appears with $|V|$ and $|E|$ of its solver graph, and the same graph under the same settings shows the same measurement in every table (Section~\ref{sec:computational-benchmark}).
\item The sample-definition stages (Table~\ref{tab:robustness-sample-size}) show how the statistic moves with the evaluation set; they are not robustness evidence, and their values are not compared with the full-cohort tables.
\end{enumerate}

\dissertationSubheading[sec:limitations-reprise]{Limitations}

The queries above were run on Arbitrum One only. Their results depend on the public log table as it stood on 9~October~2026, when scan~A and scan~B ran. The ego filter keeps only events with a cohort contract at one end, so the graphs contain no edge between two addresses outside the cohort. All reported numbers come from the machine-readable summaries of the matched cohort, the spender and trader cohorts of W0 and W1, the scaling stages and the Sybil model; the registered results of July--September come from the summary of that run (Section~\ref{sub:fresh-results}).

% !TEX root = ../main.tex
\section{Elementary Properties of PageRank and Rank Correlation}
\label{ch:appendix-theory}

Chapters~\ref{ch:literature}--\ref{ch:results} rely on the elementary properties summarized here. The material is standard: the PageRank properties follow Page et al.\cend{15} and Langville and Meyer\cend{23}, and the rank correlations Spearman\cend{20} and Kendall\cend{21}.

\dissertationSubheading[sec:pr-existence]{Existence and uniqueness of PageRank}

Take $P$ to be a row-stochastic matrix on $|V|$ nodes, possibly after the dangling-node repair of Equation~\eqref{eq:transition}. For $d\in(0,1)$ the Google matrix of Equation~\eqref{eq:google-matrix},
\[
G_{\mathrm{PR}}=d P^{\top}+\frac{1-d}{|V|}\mathbf{1}\mathbf{1}^{\top},
\]
is positive and column-stochastic. By the Perron--Frobenius theorem for positive matrices, $G_{\mathrm{PR}}$ has a single eigenvalue of modulus one, and the corresponding eigenvector $\mathbf{r}$ is strictly positive; we can scale it so that $\sum_i r_i=1$. It follows that, once the damping and the dangling-node rules are fixed, every directed graph has unique PageRank scores \cend{23}.

Two restart vectors in this thesis are not uniform. AWP's, proportional to activeness, is zero at every address that has sent nothing, and S-PR's, the normalized score of the allowance walk, is zero at every transfer-layer wallet without an allowance score, so neither Google matrix is positive. The fixed point is still unique: for two probability vectors the restart term cancels in their difference, and the operator multiplies their $\|\cdot\|_1$ distance by at most $d$. In general it is not strictly positive, however, because a wallet with no restart mass that no path from a restart wallet reaches keeps a score of zero.

Uniqueness guarantees a well-defined comparison between EndorseRank and AWP. Once its parameters are fixed, each method assigns exactly one score vector to a given edge set, so any difference between $\mathbf{s}_{\mathrm{ER}}$ and $\mathbf{s}_{\mathrm{AWP}}$ must arise from their edges and restart distributions and cannot be traced to multiple solvers.

\dissertationSubheading[sec:pr-continuity]{Continuity in edge weights}

For fixed $d\in(0,1)$, PageRank depends continuously on the entries of $P$. If $\tilde P$ is a second row-stochastic matrix on the same nodes, with PageRank vector $\tilde{\mathbf{r}}$, then
\[
\|\mathbf{r}-\tilde{\mathbf{r}}\|_1\le\frac{d}{1-d}\,\max_i\bigl\|P_{i\cdot}-\tilde P_{i\cdot}\bigr\|_1,
\]
a bound that does not depend on $|V|$ \cend{23}.

Continuity in $P$ is not the same as continuity in the edge weights. A row $P_{i\cdot}=w_{i\cdot}/o_i$ moves continuously with the weights only while the out-strength $o_i$ stays positive. When a node gains its first out-edge or loses its last one, its row jumps between the uniform dangling row and a row concentrated on its out-neighbors, however small the weight involved. In the worked example of Section~\ref{sec:pagerank-worked-example}, adding an edge $D\to A$ of weight $10^{-12}$ moves $\mathbf{r}$ by $0.37$ in the $\|\cdot\|_1$ norm and raises $r_A$ from $0.1375$ to $0.3202$.

The damping sweep is covered by continuity: it keeps $P$ fixed and changes $d$ smoothly, and for $d<1$ the PageRank vector varies continuously with $d$. The top-token filter is not covered. It removes edges, which can leave a wallet without out-edges, so the stability of that check rests on the data alone (Table~\ref{tab:robustness-tokens}). Empirically, EndorseRank's allowance family mean is unchanged in the first three decimal places across $d\in\{0.75,0.85,0.95\}$ (Table~\ref{tab:robustness-damping}); with most of the cohort tied, the ranking could hardly have moved.

\dissertationSubheading[sec:pr-sparsity]{Sparsity and per-iteration cost}

With $m=|E|$ nonzero edges, the sparse product $P^{\top}\mathbf{x}$ costs $O(m)$ arithmetic operations, and building the sparse operator from the edge list is also $O(m)$, so power iteration over $I$ iterations is $O(m+mI)$. On the matched cohort the transfer graph has about $9.3$ times as many edges as the endorsement graph, so even at identical $I$ the transfer method would run roughly an order of magnitude longer. In the measured calls the assembly of the operator takes most of the time and the iterations little (Section~\ref{sec:benchmark-results}), so the measured time ratio of $9.2$ tracks the edge ratio.

\dissertationSubheading[sec:spearman-properties]{Properties of Spearman's \texorpdfstring{$\rho$}{rho}}

Spearman's $\rho$ is Pearson's correlation computed on ranks. It lies in $[-1,1]$ and does not change under strictly increasing transformations of either variable. A value of $\rho=1$ means perfect monotone agreement, and $\rho=-1$ perfect reverse agreement. Tied values receive mid-ranks, and the evaluation pipeline relies on standard tie-aware implementations. Because reputation scores and proxies are continuous and heavy-tailed, rank-based association suits them better than Pearson correlation on the raw values. The AWP baseline paper likewise evaluates its ranking with rank statistics~\cend{3}.

\dissertationSubheading[sec:kendall-properties]{Properties of Kendall's \texorpdfstring{$\tau$}{tau}}

Kendall's $\tau$ compares the numbers of concordant and discordant pairs, $n_c$ and $n_d$. Without tie correction, $\tau_a=(n_c-n_d)/n_0$ with $n_0=n(n-1)/2$ is the probability that two randomly drawn wallets are ordered the same way by both rankings minus the probability that they are ordered oppositely \cend{21}. Every reported $\tau$ is the tie-corrected $\tau_b$ of Equation~\eqref{eq:kendall}. The two coincide without ties, and here ties are heavy. EndorseRank takes $99$ distinct values on the matched cohort, so $67.3\%$ of all wallet pairs are tied on it, against $0.5\%$ on AWP, and for EndorseRank against AWP the reported $\tau_b=0.360$ corresponds to $\tau_a=0.206$. Ties also cap $\tau_b$. EndorseRank's largest attainable value is about $0.38$ against the allowance proxies, which are positive for exactly the 131 wallets it separates, and about $0.57$ against a proxy without ties, so its $0.375$ on the allowance family (Table~\ref{tab:alignment-hybrid}) is nearly complete agreement on the pairs the rankings separate and cannot be compared directly with a coefficient that has a different ceiling.

When $n$ is large, $\tau$ and $\rho$ usually agree in sign, but they differ in magnitude. Across the pairs of EndorseRank or AWP with a proxy where $|\tau|>0.02$, the ratio $\rho/\tau$ lies between $1.01$ and $1.49$, and $\rho$ is always the larger of the two. We report both, as the AWP baseline does~\cend{3}. Our main summary statistic is the family-mean $\tau$ with a paired-bootstrap interval~\cend{34}.

\dissertationSubheading[sec:construct-validity-note]{Construct validity note}

Cronbach and Meehl~\cend{22} distinguish criterion-related validity from construct validity. Since no definitive gold standard default label exists, this research takes a construct-oriented approach. Each score is tested against validation checks that embody its conceptual definition (e.g.\ receiving allowances for EndorseRank; receiving inbound transfers for AWP), and also against a holdout period with two labels, new approval pairs and new transfer senders. High within-construct $\tau$ supports the interpretation. A low cross-construct $\tau$, by contrast, does not automatically signal failure; it may indicate discriminant validity.

\dissertationSubheading[sec:numerical-summary-sheet]{Where each primary quantity is reported}

For each primary quantity, the index below names the table where it is established.

\begin{center}
\fitwidth{%
\begin{tabular}{ll}
\toprule
Quantity & Reported in \\
\midrule
Matched-cohort size, $|V|$, $|E|$, runtime, SD, memory, iterations (four methods) & Table~\ref{tab:benchmark-runtime} \\
In-cohort scaling stages & Table~\ref{tab:benchmark-scaling-incohort} \\
Pool subsample stages ($N=99{,}080$) & Table~\ref{tab:benchmark-scaling} \\
Family-mean $\tau$ with 95\% CI, every score & Table~\ref{tab:alignment-hybrid} \\
Paired $\Delta\tau$ contrasts of C-PR and S-PR & Table~\ref{tab:tau-diff-hybrid} \\
Inter-method $\tau$; wallets outside a graph as isolated nodes & Table~\ref{tab:tie-sensitivity} \\
Damping, top-token, sample-definition sweeps & Tables~\ref{tab:robustness-damping}, \ref{tab:robustness-tokens}, \ref{tab:robustness-sample-size} \\
Spender holdout ($n=1{,}335$) & Tables~\ref{tab:holdout-spenders}, \ref{tab:holdout-diff} and~\ref{tab:holdout-receiving} \\
Exploratory trader run & Table~\ref{tab:holdout-traders} \\
Registered June--August replication & Section~\ref{sub:fresh-results} \\
Comparisons between EndorseRank and AWP, per-proxy results & Appendix~\ref{ch:appendix-er-awp} \\
Observation window & 2025-12-01 -- 2026-05-31 (Table~\ref{tab:parameter-summary}) \\
\bottomrule
\end{tabular}
}
\end{center}

Each quantity listed appears both in the evaluation summary and in a table of Chapter~\ref{ch:results} or of the appendix.

% !TEX root = ../main.tex
\section{Data Description}
\label{ch:appendix-data}

This appendix describes the extracted data behind Chapter~\ref{ch:results}: the volume of the ego extraction and the size of the two graphs, the tokens that carry the most rows and the largest amounts, how often owners grant unlimited allowances, and the activity of GMX~V2 contract accounts by quarter. The account types of the cohorts are in Table~\ref{tab:account-types}, and the extraction itself is described in Appendix~\ref{sec:bq-extraction}.

\dissertationSubheading[sec:data-events]{Extracted events and graphs}

The ego extraction of the observation window holds $389{,}866{,}882$ approval logs and $1{,}538{,}404{,}067$ transfer logs of non-zero value, each with a contract of the matched cohort at one end. From these events the allowance graph at $T_{\text{obs}}$ has $13{,}055{,}150$ nodes and $22{,}321{,}470$ edges, and the transfer graph $23{,}029{,}298$ nodes and $59{,}737{,}617$ edges. At the freeze of W0 the two graphs have $21{,}550{,}791$ and $58{,}065{,}784$ edges. The allowance graph has about a third as many edges as the transfer graph. It keeps one latest allowance per token, owner and spender and drops the revoked ones, whereas the transfer graph has an edge for every sender--recipient pair that exchanged a transfer in the window.

Membership in the cohort does not depend on either graph at $T_{\text{obs}}$. Every cohort contract received non-zero approvals in the window, but $962$ of them neither hold nor grant a positive latest allowance at $T_{\text{obs}}$, so they lie outside the allowance graph; $3{,}297$ lie outside the transfer graph, and $138$ outside both. Chapter~\ref{ch:results} scores such a contract zero under the method whose graph it misses, and Table~\ref{tab:tie-sensitivity} reports the rule that keeps it as an isolated node instead.

\dissertationSubheading[sec:data-tokens]{Tokens and allowance amounts}

The token check of Section~\ref{sub:top-tokens} rebuilds both graphs from twenty tokens, chosen once by the number of rows and once by the summed raw amount (\texttt{sql/11\_top\_tokens.sql}). Table~\ref{tab:top-tokens} lists the $27$ tokens of the two lists, which share $13$. The list by rows is led by WETH, the native USDC and USDT0, and it holds the bridged USDC, ARB, WBTC and several smaller tokens. Raw amounts are summed without regard to decimals, and an unlimited allowance counts as $2^{256}-1$ base units, so the list by amount takes in seven tokens with comparatively few rows, among them aArbWETH, the Aave deposit receipt for WETH, and two Fiat24 tokens with two decimals.

\begin{table}[htbp]
\centering
\caption[Tokens of the two top-20 lists of the token check.]{Tokens of the two top-20 lists of the token check (Table~\ref{tab:robustness-tokens}), ordered by rows: transfers in the ego tables up to $T_{\text{obs}}$ plus positive latest allowances at $T_{\text{obs}}$. By amount: among the 20 tokens with the largest summed raw amount of the same rows. By rows: among the 20 tokens with the most rows. Symbols and names are those the token contracts returned to \texttt{symbol()} and \texttt{name()} calls on 10~October~2026; the extraction keeps only addresses, which are given in full in \texttt{data/1-raw-blockchain-extracts/graph-tables/top\_tokens.csv}.}
\label{tab:top-tokens}
\fitwidth{%
\begin{tabular}{lllrcc}
\toprule
Symbol & Name & Address & Rows & By amount & By rows \\
\midrule
WETH & Wrapped Ether & \texttt{0x82af49}\ldots\texttt{3fbab1} & 509{,}394{,}850 & \checkmark & \checkmark \\
USDC & USD Coin & \texttt{0xaf88d0}\ldots\texttt{8e5831} & 366{,}294{,}721 & \checkmark & \checkmark \\
USDT0 & USDT0 & \texttt{0xfd086b}\ldots\texttt{9fcbb9} & 163{,}462{,}893 & \checkmark & \checkmark \\
ARB & Arbitrum & \texttt{0x912ce5}\ldots\texttt{9e6548} & 75{,}552{,}592 & \checkmark & \checkmark \\
WBTC & Wrapped BTC & \texttt{0x2f2a25}\ldots\texttt{fc5b0f} & 71{,}824{,}604 & \checkmark & \checkmark \\
USDC & USD Coin (Arb1) & \texttt{0xff970a}\ldots\texttt{db5cc8} & 71{,}329{,}831 & \checkmark & \checkmark \\
PENDLE & Pendle & \texttt{0x0c880f}\ldots\texttt{a8c9e8} & 19{,}516{,}891 & \checkmark & \checkmark \\
WBTC\_ETH\_FS\_LT & WBTC ETH Futureswap Liquidity Token & \texttt{0xa522a5}\ldots\texttt{0d8d98} & 17{,}735{,}463 &  & \checkmark \\
NPC & NPC & \texttt{0xe1a498}\ldots\texttt{0bb8b0} & 12{,}163{,}764 &  & \checkmark \\
UXUY & UXUY Token & \texttt{0xe2035f}\ldots\texttt{7c79f9} & 9{,}761{,}886 & \checkmark & \checkmark \\
AIUS & Arbius & \texttt{0x4a24b1}\ldots\texttt{0cabc3} & 8{,}804{,}452 &  & \checkmark \\
DAI & Dai Stablecoin & \texttt{0xda1000}\ldots\texttt{000da1} & 8{,}677{,}480 & \checkmark & \checkmark \\
LINK & ChainLink Token & \texttt{0xf97f4d}\ldots\texttt{539fb4} & 8{,}675{,}046 &  & \checkmark \\
xUSD & xUSD & \texttt{0xe80772}\ldots\texttt{afca65} & 6{,}437{,}746 &  & \checkmark \\
ZRO & LayerZero & \texttt{0x698588}\ldots\texttt{3271cd} & 6{,}318{,}171 & \checkmark & \checkmark \\
UXLINK & UXLINK Token & \texttt{0x1a6b3a}\ldots\texttt{c8cff1} & 6{,}187{,}243 &  & \checkmark \\
VSN & Vision & \texttt{0x6fbbbd}\ldots\texttt{7db74b} & 6{,}052{,}221 & \checkmark & \checkmark \\
wstETH & Wrapped liquid staked Ether 2.0 & \texttt{0x5979d7}\ldots\texttt{800529} & 5{,}553{,}387 &  & \checkmark \\
GRT & Graph Token & \texttt{0x962306}\ldots\texttt{7e88c7} & 4{,}761{,}427 & \checkmark & \checkmark \\
GMX & GMX & \texttt{0xfc5a1a}\ldots\texttt{35ad0a} & 4{,}378{,}643 & \checkmark & \checkmark \\
BEBE & BEBE & \texttt{0x406d59}\ldots\texttt{64f97b} & 2{,}720{,}903 & \checkmark &  \\
STG & StargateToken & \texttt{0x669434}\ldots\texttt{ce1eb6} & 1{,}647{,}816 & \checkmark &  \\
aArbWETH & Aave Arbitrum WETH & \texttt{0xe50fa9}\ldots\texttt{8128c8} & 1{,}343{,}622 & \checkmark &  \\
OTMT & OTMAirdrop & \texttt{0x6b708d}\ldots\texttt{37f770} & 493{,}833 & \checkmark &  \\
BPET & BPET & \texttt{0x6daf58}\ldots\texttt{862ac5} & 389{,}567 & \checkmark &  \\
USD24 & Fiat24 USD & \texttt{0xbe00f3}\ldots\texttt{cc0d62} & 207{,}237 & \checkmark &  \\
CNH24 & Fiat24 CNH & \texttt{0x7288ac}\ldots\texttt{0144fd} & 102{,}957 & \checkmark &  \\
\bottomrule
\end{tabular}
}
\end{table}

Unlimited allowances are common among the owners of the allowance graph. Of the $13{,}029{,}306$ owners with a positive latest allowance at $T_{\text{obs}}$, $5{,}798{,}712$ hold at least one allowance of $1.15\times10^{77}$ base units or more, close to the largest \texttt{uint256} value $2^{256}-1\approx1.16\times10^{77}$, and $1{,}532{,}301$ hold both such an allowance and a smaller one (\texttt{sql/10\_describe.sql}). EndorseRank and AWP pass amounts through the bounded value weight, under which an unlimited allowance counts about the same as any other allowance of ten base units or more. The raw allowance layer of C-PR sums amounts without a bound, so in the row of an owner with both kinds the unlimited allowances take nearly all the weight unless the other allowances are of a similar size (Section~\ref{sub:coupled-graph}).

\dissertationSubheading[sec:data-gmx]{GMX~V2 contract traders by quarter}

The trader cohorts are GMX~V2 accounts that hold contract code (Section~\ref{sub:cohorts}). Table~\ref{tab:gmx-quarters} counts their position closes, the \texttt{PositionDecrease} events of the protocol's event emitter \parencite{gmxnddocs}, by calendar quarter, together with the liquidations among them and the number of contract accounts that closed at least one and at least three positions in the quarter.

\begin{table}[htbp]
\centering
\caption[GMX~V2 closes of contract accounts by quarter.]{GMX~V2 closes of contract accounts by quarter. October~2023 to June~2026: closes by the $7{,}633$ contracts among the $40{,}307$ GMX~V2 accounts with at least three closes in the observation window (Table~\ref{tab:account-types}). July to September~2026 (the label window of W1): closes by the contract accounts of the observation window and of the W1 run. Liquidation: a close with the liquidation order type or emitted through the liquidation handler. Source: the decoded closes in \texttt{data/2-processed-tables-and-evaluations/gmx/}.}
\label{tab:gmx-quarters}
\fitwidth{%
\begin{tabular}{lrrrr}
\toprule
Quarter & Closes & Liquidations & Accounts with a close & Accounts with $\geq 3$ closes \\
\midrule
2023 Q4 & 1{,}119 & 104 & 180 & 105 \\
2024 Q1 & 7{,}956 & 692 & 1{,}247 & 961 \\
2024 Q2 & 11{,}509 & 589 & 2{,}120 & 1{,}165 \\
2024 Q3 & 7{,}947 & 456 & 1{,}807 & 724 \\
2024 Q4 & 7{,}343 & 904 & 1{,}321 & 698 \\
2025 Q1 & 6{,}929 & 498 & 497 & 273 \\
2025 Q2 & 12{,}080 & 379 & 368 & 250 \\
2025 Q3 & 28{,}753 & 897 & 1{,}750 & 1{,}567 \\
2025 Q4 & 40{,}240 & 1{,}596 & 1{,}915 & 1{,}744 \\
2026 Q1 & 6{,}026 & 233 & 269 & 171 \\
2026 Q2 & 3{,}783 & 125 & 169 & 127 \\
\midrule
2026 Q3 (W1) & 6{,}805 & 182 & 165 & 131 \\
\midrule
October~2023 -- June~2026 & 133{,}685 & 6{,}473 & & \\
\bottomrule
\end{tabular}
}
\end{table}

Contract accounts closed the most positions in the second half of 2025, $28{,}753$ and $40{,}240$ in the third and fourth quarters, when $1{,}567$ and $1{,}744$ of them closed at least three. In the first half of 2026 their closes fell to $6{,}026$ and $3{,}783$, and the accounts with at least three closes to $171$ and $127$. A trader cohort further requires the account to be a node of the graph at the freeze: of the $127$ contract accounts with at least three closes in April--June~2026, $88$ form the trader cohort of W0, and of the $131$ in July--September, $70$ form that of W1. These small cohorts are why every trader result of Chapter~\ref{ch:results} has wide intervals.

% !TEX root = ../main.tex
\section{EndorseRank and AWP Side by Side}
\label{ch:appendix-er-awp}

Chapter~\ref{ch:results} reports every score of the thesis together and reads each PageRank against the raw degree it smooths. This appendix collects the direct comparisons between the two single-layer scores, EndorseRank and AWP: their family means and paired differences in the same window, their per-proxy coefficients, their paired differences out of window, and their contrasts on every trader label. None of the comparisons on the matched cohort or on the spender cohorts was fixed in advance, and none of them carries a claim of the thesis; the trader pairs were registered with W1 (Section~\ref{sec:er-awp-neutral-grid}). The two scores walk on different edges and are built for different constructs, so on each check over a full cohort the score ahead is the one whose edges the check counts. Outcomes built from neither edge type are the exception, since neither score was built for them; Section~\ref{sub:neutral-label-results} reports them.

Coverage differs between the two. Of the matched cohort, $962$ contracts lie outside the allowance graph and $3{,}297$ outside the transfer graph, and $3{,}753$ of the $33{,}101$ spenders of W0 had no transfer edge before the freeze, so AWP scores them zero. EndorseRank, by contrast, leaves almost no pair of cohort contracts tied (Section~\ref{sec:rank-divergence}). Out of window the two scores stand differently against the raw degrees they smooth: AWP ranks later new senders above the transfer in-degree in both windows, while EndorseRank ranks later new approval pairs below the in-approve degree in both (Section~\ref{sec:holdout-results}). A lead of EndorseRank over AWP therefore does not show that propagation over approvals adds anything to the count of approving owners.

\dissertationSubheading[sec:er-awp-same-window]{Same window}

Table~\ref{tab:alignment-hybrid} gives the family means of the two scores with their intervals, and the last row of Table~\ref{tab:tie-sensitivity} their agreement. The two patterns cross. AWP tracks the transfer family most closely ($0.447$, interval $[0.440,0.453]$) and the allowance family least ($0.173$, $[0.165,0.180]$), while EndorseRank tracks the allowance family most closely ($0.481$, $[0.476,0.487]$) and stands at $0.122$ ($[0.113,0.129]$) on the transfer family. Table~\ref{tab:tau-diff} tests five paired differences on the same contract resamples (Section~\ref{sec:statistical-tests}), numbered i to v in the order of its rows.

% Generated by experiments/scripts/export_latex.py -- do not edit by hand.
% Source: experiments/data/2-processed-tables-and-evaluations/same-window/eval_summary.json
\begin{table}[htbp]
\centering
\caption[Same-window paired contrasts between EndorseRank and AWP.]{Paired bootstrap tests of differences between Kendall $\tau$ coefficients of EndorseRank and AWP in the same window (matched cohort, $n=27{,}844$ contracts). Family entries use the family-mean $\tau$; the same contract resamples are used for $a$ and $b$, so the interval is on the paired difference. An interval excluding 0 indicates a difference not attributable to sampling variation at the 5\% level. CI = 95\% percentile bootstrap over contracts (400 paired resamples, seed 42).}
\label{tab:tau-diff}
\fitwidth{%
\begin{tabular}{lccccc}
\toprule
Contrast ($a$ minus $b$) & $\tau_a$ & $\tau_b$ & $\Delta\tau$ & 95\% CI of $\Delta\tau$ & $P(\Delta\tau>0)$ \\
\midrule
EndorseRank allowance minus EndorseRank transfer & 0.481 & 0.122 & +0.359 & [0.351, 0.367] & 1.000 \\
AWP transfer minus EndorseRank transfer & 0.447 & 0.122 & +0.325 & [0.315, 0.335] & 1.000 \\
EndorseRank allowance minus AWP allowance & 0.481 & 0.173 & +0.309 & [0.300, 0.317] & 1.000 \\
AWP sybil-stability minus EndorseRank sybil-stability & 0.311 & 0.025 & +0.285 & [0.277, 0.294] & 1.000 \\
EndorseRank in-approve degree minus EndorseRank in-degree & 0.626 & 0.196 & +0.430 & [0.421, 0.440] & 1.000 \\
\bottomrule
\end{tabular}
}
\end{table}


All five intervals lie above zero. EndorseRank's allowance coefficient exceeds its transfer coefficient (contrast~i, $+0.359$, interval $[0.351,0.367]$); AWP's transfer coefficient exceeds EndorseRank's (contrast~ii, $+0.325$, $[0.315,0.335]$); EndorseRank's allowance coefficient exceeds AWP's (contrast~iii, $+0.309$, $[0.300,0.317]$); AWP leads on the Sybil-stability family (contrast~iv, $+0.285$, $[0.277,0.294]$); and EndorseRank follows the in-approve degree more closely than the in-degree (contrast~v, $+0.430$, $[0.421,0.440]$). Contrasts~i and~v compare EndorseRank with itself across proxies. EndorseRank ties only $0.15\%$ of the contract pairs, so its own ties cap none of these coefficients, and the ties of the proxies do not explain gaps of this size: on the in-degree AWP reaches $0.541$ ($[0.534,0.547]$), where EndorseRank reaches $0.196$ ($[0.187,0.204]$). Keeping the contracts outside the allowance graph as isolated nodes changes none of the family means (Table~\ref{tab:tie-sensitivity}).

The opposite signs of contrasts~ii and~iii follow from the constructs. Each score is ahead on the family built from its own edges, and each of these families includes the in-degree that the corresponding PageRank smooths. The agreement is partly built in and checks the intended construct of each score; it is not evidence from outside the score. Every contract of the cohort received approvals from at least three owners, and AWP's allowance coefficients lie above zero (Table~\ref{tab:alignment-allowance}): contracts that many owners approve also tend to be central in the transfer graph. Neither difference says which score is the better measure of reputation.

On the scaling stages of Table~\ref{tab:robustness-sample-size}, whose values are not comparable with those of the full cohort, the order of the two scores within each family is the same at every stage: AWP is ahead on the transfer and Sybil-stability families and EndorseRank on the allowance family.

\dissertationSubheading[sec:proxy-family-detail]{Proxy detail}

Each family mean condenses several per-proxy coefficients. Tables~\ref{tab:alignment-transfer} to~\ref{tab:alignment-sybil-stability} report them one proxy at a time for EndorseRank and AWP, as Spearman $\rho$, Kendall $\tau$ and the 95\% interval of $\tau$. Before correlation, every proxy is oriented so that a higher value means higher expected reputation, and the transfer proxies are the ones used in the evaluation of AWP by its authors \parencite{do2023awp}.

\dissertationSubsubheading[sub:transfer-family]{Transfer family}

Table~\ref{tab:alignment-transfer} gives the coefficients of the two scores with the in-degree and the in-value.

% Generated by experiments/scripts/export_latex.py -- do not edit by hand.
% Source: experiments/data/2-processed-tables-and-evaluations/same-window/eval_summary.json
\begin{table}[htbp]
\centering
\caption[Domain alignment of EndorseRank and AWP with the transfer proxies.]{Domain alignment of EndorseRank and AWP with the transfer proxies (in-degree and in-value of each contract), same window (matched cohort, $n=27{,}844$ contracts). Kendall $\tau$ and Spearman $\rho$ between scores and proxy rankings. CI = 95\% percentile bootstrap over contracts (400 paired resamples, seed 42).}
\label{tab:alignment-transfer}
\fitwidth{%
\begin{tabular}{lccc}
\toprule
Proxy & Spearman $\rho$ & Kendall $\tau$ & 95\% CI of $\tau$ \\
\midrule
\multicolumn{4}{l}{\textit{EndorseRank}} \\
In-degree & 0.283 & 0.196 & [0.187, 0.204] \\
In-value (sum inbound) & 0.081 & 0.048 & [0.039, 0.055] \\
\addlinespace
\multicolumn{4}{l}{\textit{AWP}} \\
In-degree & 0.706 & 0.541 & [0.534, 0.547] \\
In-value (sum inbound) & 0.496 & 0.353 & [0.345, 0.361] \\
\bottomrule
\end{tabular}
}
\end{table}


AWP agrees with the in-degree at $0.541$ (interval $[0.534,0.547]$) and with the in-value at $0.353$ ($[0.345,0.361]$). AWP's value weight saturates near ten base units (Section~\ref{sub:awp-graph}), so the size of a transfer barely enters its score, while the in-value sums raw amounts. EndorseRank's transfer coefficients lie above zero, at $0.196$ ($[0.187,0.204]$) and $0.048$ ($[0.039,0.055]$), and their intervals lie far below AWP's: contracts that many owners approve receive transfers from somewhat more senders than the rest of the cohort. Out of window AWP ranks later new senders above the raw in-degree in both windows (Section~\ref{sec:holdout-results}), so on the transfer side the walk adds to the count it smooths.

\dissertationSubsubheading[sub:allowance-family]{Allowance family}

Table~\ref{tab:alignment-allowance} gives the coefficients of the two scores with the in-approve degree and the in-approve value.

% Generated by experiments/scripts/export_latex.py -- do not edit by hand.
% Source: experiments/data/2-processed-tables-and-evaluations/revised-usd/same-window/eval_summary.json
\begin{table}[htbp]
\centering
\caption[Domain alignment of EndorseRank and AWP with the allowance proxies.]{Domain alignment of EndorseRank and AWP with the allowance proxies (in-approve degree and value of each spender contract), same window (matched cohort, $n=15{,}107$ contracts). Kendall $\tau$ and Spearman $\rho$ between scores and proxy rankings. CI = 95\% percentile bootstrap over contracts (400 paired resamples, seed 42).}
\label{tab:alignment-allowance}
\fitwidth{%
\begin{tabular}{lccc}
\toprule
Proxy & Spearman $\rho$ & Kendall $\tau$ & 95\% CI of $\tau$ \\
\midrule
\multicolumn{4}{l}{\textit{EndorseRank}} \\
In-approve degree & 0.811 & 0.655 & [0.648, 0.662] \\
In-approve value (sum) & 0.681 & 0.515 & [0.506, 0.524] \\
\addlinespace
\multicolumn{4}{l}{\textit{AWP}} \\
In-approve degree & 0.286 & 0.206 & [0.193, 0.217] \\
In-approve value (sum) & 0.260 & 0.180 & [0.169, 0.191] \\
\bottomrule
\end{tabular}
}
\end{table}


EndorseRank agrees with the in-approve degree at $0.626$ (interval $[0.620,0.632]$) and with the in-approve value at $0.336$ ($[0.329,0.344]$). Under the bounded value weight nearly every allowance counts the same, so EndorseRank follows the number of approving owners more than the amounts they approve. For a spender whose owners receive no approvals it equals the closed form of Section~\ref{sub:theory-closed-form}, a function of the owner-normalized in-strength. AWP agrees with the two proxies at $0.227$ ($[0.218,0.235]$) and $0.118$ ($[0.110,0.126]$). The comparison is partly built in, since EndorseRank is computed from the same approvals as the proxies.

\dissertationSubsubheading[sub:sybil-family]{Sybil-adjusted stability family}

Table~\ref{tab:alignment-sybil-stability} gives the coefficients of the two scores with the inbound counterparty ratio, the transfer tenure and the active months.

% Generated by experiments/scripts/export_latex.py -- do not edit by hand.
% Source: experiments/data/2-processed-tables-and-evaluations/revised-usd/same-window/eval_summary.json
\begin{table}[htbp]
\centering
\caption[Domain alignment of EndorseRank and AWP with the Sybil-adjusted stability proxies.]{Domain alignment of EndorseRank and AWP with the Sybil-adjusted stability proxies from transfer activity, same window (matched cohort, $n=15{,}107$ contracts). Kendall $\tau$ and Spearman $\rho$ between scores and proxy rankings. CI = 95\% percentile bootstrap over contracts (400 paired resamples, seed 42).}
\label{tab:alignment-sybil-stability}
\fitwidth{%
\begin{tabular}{lccc}
\toprule
Proxy & Spearman $\rho$ & Kendall $\tau$ & 95\% CI of $\tau$ \\
\midrule
\multicolumn{4}{l}{\textit{EndorseRank}} \\
Inbound counterparty ratio & -0.231 & -0.156 & [-0.166, -0.147] \\
Transfer tenure (days) & 0.124 & 0.080 & [0.069, 0.091] \\
Active months & 0.146 & 0.098 & [0.087, 0.109] \\
\addlinespace
\multicolumn{4}{l}{\textit{AWP}} \\
Inbound counterparty ratio & 0.234 & 0.127 & [0.114, 0.141] \\
Transfer tenure (days) & 0.683 & 0.508 & [0.499, 0.517] \\
Active months & 0.704 & 0.549 & [0.540, 0.558] \\
\bottomrule
\end{tabular}
}
\end{table}


The table shows AWP ahead of EndorseRank on the transfer tenure ($0.452$, interval $[0.445,0.458]$, against $0.083$, $[0.074,0.090]$) and on the active months ($0.499$, $[0.492,0.506]$, against $0.121$, $[0.113,0.129]$), as one would expect from a transfer score that rewards activity spread over the window; the paired difference on the family mean excludes zero (Table~\ref{tab:tau-diff}, contrast~iv). On the inbound counterparty ratio both coefficients are negative, $-0.128$ ($[-0.135,-0.121]$) for EndorseRank and $-0.019$ ($[-0.029,-0.010]$) for AWP. The ratio divides the distinct other senders by the transfers received from them, so it falls when the same senders return; EndorseRank's negative coefficient says that contracts that many owners approve are, on this measure, contracts whose senders return. The family describes how much longitudinal transfer information each score carries.

\dissertationSubheading[sec:er-awp-holdout]{Out of window}

Table~\ref{tab:endorserank-awp-holdout} collects the paired differences between the two scores in both holdout windows. In every row $a$ is EndorseRank and $b$ AWP.

% Generated by experiments/scripts/export_latex.py -- do not edit by hand.
% Source: experiments/data/2-processed-tables-and-evaluations/revised-usd/holdout-w0/eval_summary.json
% Source: experiments/data/2-processed-tables-and-evaluations/revised-usd/supplementary/supplementary.json
% Source: experiments/data/2-processed-tables-and-evaluations/revised-usd/window-w1/eval_summary.json
\begin{table}[htbp]
\centering
\caption[Paired bootstrap differences between EndorseRank and AWP out of window.]{Paired bootstrap differences between EndorseRank and AWP out of window (W0 spenders $n=17{,}050$; W0 traders $n=88$; W1 spenders $n=18{,}152$; W1 traders $n=70$). None of them enters rule A or rule B: the spender rows were computed after the labels were known, and the trader rows are the EndorseRank-minus-AWP pair of the comparison on an outcome built from neither edge type (Tables~\ref{tab:neutral-label-contrasts} and~\ref{tab:neutral-label-grid}). The W1 spender rows were computed from the saved W1 scores after the W1 evaluation, with the resamples of Table~\ref{tab:fresh-holdout}. Within each block $a$ and $b$ share the contract resamples (400 for the spenders and 2{,}000 for the traders; seed 42).}
\label{tab:endorserank-awp-holdout}
\fitwidth{%
\begin{tabular}{llccccc}
\toprule
Contrast ($a$ minus $b$) & Label & $\tau_a$ & $\tau_b$ & $\Delta\tau$ & 95\% CI of $\Delta\tau$ & $P(\Delta\tau>0)$ \\
\midrule
\multicolumn{7}{l}{\emph{W0, spenders ($n=17{,}050$)}} \\
EndorseRank minus AWP & new approval pairs & 0.231 & 0.196 & +0.035 & [0.019, 0.051] & 1.000 \\
EndorseRank minus AWP & new transfer senders & 0.200 & 0.346 & -0.146 & [-0.156, -0.136] & 0.000 \\
\midrule
\multicolumn{7}{l}{\emph{W0, traders with at least three closes ($n=88$)}} \\
EndorseRank minus AWP & liquidation-free close rate & 0.224 & 0.240 & -0.017 & [-0.156, 0.128] & 0.394 \\
\midrule
\multicolumn{7}{l}{\emph{W1, spenders ($n=18{,}152$)}} \\
EndorseRank minus AWP & new approval pairs & 0.234 & 0.169 & +0.065 & [0.051, 0.081] & 1.000 \\
EndorseRank minus AWP & new transfer senders & 0.194 & 0.330 & -0.136 & [-0.146, -0.127] & 0.000 \\
\midrule
\multicolumn{7}{l}{\emph{W1, traders with at least three closes ($n=70$)}} \\
EndorseRank minus AWP & liquidation-free close rate & 0.297 & 0.329 & -0.032 & [-0.133, 0.066] & 0.279 \\
\bottomrule
\end{tabular}
}
\end{table}


On the spenders of W0, EndorseRank ranks new approval pairs at $\tau=0.241$ (interval $[0.234,0.248]$) against $0.196$ ($[0.187,0.206]$) for AWP, a paired difference of $+0.045$ ($[0.033,0.055]$). On new transfer senders the order reverses: AWP reaches $0.340$ ($[0.333,0.348]$) against $0.226$ ($[0.219,0.234]$) for EndorseRank, a difference of $-0.114$ ($[-0.120,-0.108]$). Every label falls after all scored events, so on this cohort each score ranks the later arrivals of its own edge type better than the other score does. Coverage accounts for part of the first gap. AWP scores zero for the $3{,}753$ spenders without a transfer edge before the freeze, and $11.9\%$ of them gained a new approval pair, against $10.1\%$ of the rest. On the $29{,}228$ spenders that had received a transfer from another address, the order on new approval pairs reverses: AWP reaches $0.262$ ($[0.253,0.270]$) and EndorseRank $0.249$ ($[0.241,0.256]$), without a paired test (Table~\ref{tab:holdout-receiving}). On new senders AWP stays ahead there, at $0.349$ ($[0.342,0.356]$) against $0.247$ ($[0.240,0.255]$).

In July--September, with the contrasts computed from the saved scores of W1 after the registered evaluation, the pattern on the full spender cohort repeats. EndorseRank ranks new approval pairs at $0.233$ ($[0.226,0.240]$) against $0.177$ ($[0.166,0.187]$) for AWP ($+0.056$, $[0.046,0.068]$), and AWP ranks new senders at $0.322$ ($[0.316,0.329]$) against $0.210$ ($[0.203,0.217]$) ($-0.112$, $[-0.119,-0.106]$). On the liquidation-free close rate of the traders, one of the pairs registered with W1, EndorseRank reaches $0.230$ ($[0.115,0.341]$) in W0 and $0.306$ ($[0.177,0.421]$) in W1, and AWP $0.168$ ($[0.022,0.315]$) and $0.240$ ($[0.094,0.365]$). Neither difference can be told from zero ($+0.062$, $[-0.079,0.217]$, and $+0.066$, $[-0.069,0.227]$); the trader cohorts hold $88$ and $70$ contracts, so the intervals are wide.

On the W0 traders' profit labels (Table~\ref{tab:holdout-traders}), every interval of either score on profitable closes and realized gain includes zero: EndorseRank reaches $0.156$ ($[-0.022,0.302]$) and $0.088$ ($[-0.079,0.242]$), AWP $0.138$ ($[-0.020,0.281]$) and $0.134$ ($[-0.018,0.283]$). On the two shares both scores are negative. AWP's intervals lie below zero on both, at $-0.170$ ($[-0.299,-0.032]$) on the profitable share and $-0.253$ ($[-0.384,-0.112]$) on the non-loss share; EndorseRank's lies below zero on the non-loss share only, at $-0.170$ ($[-0.286,-0.044]$), against $-0.053$ ($[-0.186,0.082]$) on the profitable share. The transfer in-degree is negative on both shares as well, at $-0.198$ ($[-0.332,-0.053]$) and $-0.190$ ($[-0.344,-0.036]$). Neither score ranks these contract traders by trading success.

\dissertationSubheading[sec:er-awp-neutral-grid]{All trader labels in both windows}

Table~\ref{tab:neutral-label-grid} sets the allowance side against the transfer side on all six trader labels in both windows. Its four pairs and six labels were registered with W1, before the logs of July--September were extracted (Section~\ref{sub:neutral-label-protocol}). By then the single-score coefficients of the W0 traders on the profit labels were known (Table~\ref{tab:holdout-traders}), while no score had been computed against their liquidation labels. Only P and Q on the liquidation-free close rate of W1 enter a decision, through rules R1 and R2 (Claim~C6, Section~\ref{sub:neutral-label-results}); the column for EndorseRank minus AWP is reported beside them.

% Generated by experiments/scripts/export_latex.py -- do not edit by hand.
% Source: experiments/data/2-processed-tables-and-evaluations/revised-usd/window-w1/eval_summary.json
\begin{table}[htbp]
\centering
\caption[Allowance-side minus transfer-side scores on every trader label in both windows.]{Allowance-side minus transfer-side scores, $\Delta\tau$ with 95\% paired bootstrap interval, on all six trader labels of the window W1 ($n=70$ GMX~V2 contract traders) and of W0 ($n=88$). P: in-approve degree minus transfer in-degree. Q: EndorseRank with activity restarts minus AWP (same edge weights and restart rule). C-PR: $\lambda=1$ minus $\lambda=0$ (same USD build). ER: EndorseRank minus AWP. Profitable closes and realized gain grow with the number of closes; the other four labels do not. CI = 95\% percentile bootstrap over contracts (2{,}000 paired resamples, seed 42).}
\label{tab:neutral-label-grid}
\fitwidth{%
\begin{tabular}{lcccc}
\toprule
Label & P (counts) & Q (PageRanks) & C-PR layers & ER minus AWP \\
\midrule
\multicolumn{5}{l}{\emph{Window W1, July--September 2026 ($n=70$ traders)}} \\
Liquidation-free close rate & \begin{tabular}[c]{@{}c@{}}0.084\\[-1pt]{\footnotesize [-0.013, 0.176]}\end{tabular} & \begin{tabular}[c]{@{}c@{}}-0.041\\[-1pt]{\footnotesize [-0.135, 0.050]}\end{tabular} & \begin{tabular}[c]{@{}c@{}}-0.049\\[-1pt]{\footnotesize [-0.169, 0.060]}\end{tabular} & \begin{tabular}[c]{@{}c@{}}-0.032\\[-1pt]{\footnotesize [-0.133, 0.066]}\end{tabular} \\
No liquidation & \begin{tabular}[c]{@{}c@{}}0.078\\[-1pt]{\footnotesize [-0.012, 0.166]}\end{tabular} & \begin{tabular}[c]{@{}c@{}}-0.039\\[-1pt]{\footnotesize [-0.138, 0.062]}\end{tabular} & \begin{tabular}[c]{@{}c@{}}-0.051\\[-1pt]{\footnotesize [-0.170, 0.064]}\end{tabular} & \begin{tabular}[c]{@{}c@{}}-0.029\\[-1pt]{\footnotesize [-0.136, 0.080]}\end{tabular} \\
Profitable closes & \begin{tabular}[c]{@{}c@{}}-0.143\\[-1pt]{\footnotesize [-0.318, 0.013]}\end{tabular} & \begin{tabular}[c]{@{}c@{}}-0.231\\[-1pt]{\footnotesize [-0.429, -0.027]}\end{tabular} & \begin{tabular}[c]{@{}c@{}}-0.190\\[-1pt]{\footnotesize [-0.411, 0.037]}\end{tabular} & \begin{tabular}[c]{@{}c@{}}-0.214\\[-1pt]{\footnotesize [-0.423, -0.012]}\end{tabular} \\
Realized gain & \begin{tabular}[c]{@{}c@{}}-0.135\\[-1pt]{\footnotesize [-0.308, 0.025]}\end{tabular} & \begin{tabular}[c]{@{}c@{}}-0.350\\[-1pt]{\footnotesize [-0.536, -0.157]}\end{tabular} & \begin{tabular}[c]{@{}c@{}}-0.404\\[-1pt]{\footnotesize [-0.624, -0.182]}\end{tabular} & \begin{tabular}[c]{@{}c@{}}-0.346\\[-1pt]{\footnotesize [-0.547, -0.144]}\end{tabular} \\
Profitable share & \begin{tabular}[c]{@{}c@{}}-0.196\\[-1pt]{\footnotesize [-0.364, -0.046]}\end{tabular} & \begin{tabular}[c]{@{}c@{}}-0.336\\[-1pt]{\footnotesize [-0.523, -0.141]}\end{tabular} & \begin{tabular}[c]{@{}c@{}}-0.344\\[-1pt]{\footnotesize [-0.551, -0.120]}\end{tabular} & \begin{tabular}[c]{@{}c@{}}-0.316\\[-1pt]{\footnotesize [-0.517, -0.112]}\end{tabular} \\
Non-loss share & \begin{tabular}[c]{@{}c@{}}-0.252\\[-1pt]{\footnotesize [-0.391, -0.120]}\end{tabular} & \begin{tabular}[c]{@{}c@{}}-0.350\\[-1pt]{\footnotesize [-0.526, -0.175]}\end{tabular} & \begin{tabular}[c]{@{}c@{}}-0.324\\[-1pt]{\footnotesize [-0.531, -0.109]}\end{tabular} & \begin{tabular}[c]{@{}c@{}}-0.346\\[-1pt]{\footnotesize [-0.526, -0.161]}\end{tabular} \\
\midrule
\multicolumn{5}{l}{\emph{Window W0, April--June 2026 ($n=88$ traders)}} \\
Liquidation-free close rate & \begin{tabular}[c]{@{}c@{}}0.182\\[-1pt]{\footnotesize [0.052, 0.314]}\end{tabular} & \begin{tabular}[c]{@{}c@{}}-0.014\\[-1pt]{\footnotesize [-0.147, 0.130]}\end{tabular} & \begin{tabular}[c]{@{}c@{}}-0.041\\[-1pt]{\footnotesize [-0.187, 0.105]}\end{tabular} & \begin{tabular}[c]{@{}c@{}}-0.017\\[-1pt]{\footnotesize [-0.156, 0.128]}\end{tabular} \\
No liquidation & \begin{tabular}[c]{@{}c@{}}0.197\\[-1pt]{\footnotesize [0.059, 0.333]}\end{tabular} & \begin{tabular}[c]{@{}c@{}}-0.010\\[-1pt]{\footnotesize [-0.152, 0.146]}\end{tabular} & \begin{tabular}[c]{@{}c@{}}-0.051\\[-1pt]{\footnotesize [-0.198, 0.102]}\end{tabular} & \begin{tabular}[c]{@{}c@{}}-0.012\\[-1pt]{\footnotesize [-0.154, 0.137]}\end{tabular} \\
Profitable closes & \begin{tabular}[c]{@{}c@{}}0.075\\[-1pt]{\footnotesize [-0.080, 0.222]}\end{tabular} & \begin{tabular}[c]{@{}c@{}}-0.016\\[-1pt]{\footnotesize [-0.150, 0.118]}\end{tabular} & \begin{tabular}[c]{@{}c@{}}-0.067\\[-1pt]{\footnotesize [-0.222, 0.087]}\end{tabular} & \begin{tabular}[c]{@{}c@{}}-0.005\\[-1pt]{\footnotesize [-0.163, 0.155]}\end{tabular} \\
Realized gain & \begin{tabular}[c]{@{}c@{}}0.017\\[-1pt]{\footnotesize [-0.150, 0.176]}\end{tabular} & \begin{tabular}[c]{@{}c@{}}-0.107\\[-1pt]{\footnotesize [-0.252, 0.044]}\end{tabular} & \begin{tabular}[c]{@{}c@{}}-0.248\\[-1pt]{\footnotesize [-0.433, -0.067]}\end{tabular} & \begin{tabular}[c]{@{}c@{}}-0.118\\[-1pt]{\footnotesize [-0.286, 0.046]}\end{tabular} \\
Profitable share & \begin{tabular}[c]{@{}c@{}}0.045\\[-1pt]{\footnotesize [-0.113, 0.184]}\end{tabular} & \begin{tabular}[c]{@{}c@{}}0.046\\[-1pt]{\footnotesize [-0.071, 0.172]}\end{tabular} & \begin{tabular}[c]{@{}c@{}}0.027\\[-1pt]{\footnotesize [-0.112, 0.181]}\end{tabular} & \begin{tabular}[c]{@{}c@{}}0.069\\[-1pt]{\footnotesize [-0.060, 0.208]}\end{tabular} \\
Non-loss share & \begin{tabular}[c]{@{}c@{}}-0.080\\[-1pt]{\footnotesize [-0.228, 0.058]}\end{tabular} & \begin{tabular}[c]{@{}c@{}}0.055\\[-1pt]{\footnotesize [-0.070, 0.192]}\end{tabular} & \begin{tabular}[c]{@{}c@{}}0.037\\[-1pt]{\footnotesize [-0.111, 0.189]}\end{tabular} & \begin{tabular}[c]{@{}c@{}}0.076\\[-1pt]{\footnotesize [-0.063, 0.223]}\end{tabular} \\
\bottomrule
\end{tabular}
}
\end{table}


In W1 no pair separates the two sides on either liquidation label: every interval includes zero, P on the liquidation-free close rate at $+0.063$ ($[-0.026,0.148]$) and the EndorseRank pair at $+0.066$ ($[-0.069,0.227]$). In W0 the counts favor the allowance side on both liquidation labels, P at $+0.252$ ($[0.093,0.413]$) on the liquidation-free close rate and $+0.270$ ($[0.098,0.441]$) on the no-liquidation flag, while the three PageRank pairs lean the same way with intervals that include zero. On the four profit labels of W1 every pair is negative. The intervals of Q and of the EndorseRank pair lie below zero on all four, those of P on the two shares, and those of the C-PR layers on realized gain and on both shares. In W0 most intervals on the profit labels include zero, and Q and the EndorseRank pair lean to the allowance side on the profitable share ($+0.093$, $[0.001,0.187]$, and $+0.117$, $[0.004,0.227]$), so the profit labels show no direction that holds in both windows. With $70$ and $88$ traders, none of these differences is precise.

% !TEX root = ../main.tex
\section{Claim-to-Evidence Map}
\label{ch:appendix-claims}

This appendix maps each hypothesis, research question, and empirical claim to the relevant evidence in Chapter~\ref{ch:results}.

\dissertationSubheading[sec:map-h1]{H1 / RQ1: Rank divergence}

\begin{itemize}
\item Statement: on the same cohort, EndorseRank and AWP rank wallets differently.
\item Evidence: Section~\ref{sec:rank-divergence}; Table~\ref{tab:tie-sensitivity}, whose last row gives the inter-method $\tau_b$ with its interval, also with the wallets outside each graph kept as isolated nodes; Figures~\ref{fig:rank-scatter} and~\ref{fig:score-distributions}.
\item Interpretation: the rankings are not equivalent, and their agreement comes from the wallets that EndorseRank scores zero; the O1 criterion is met (Table~\ref{tab:objectives-achievement}).
\end{itemize}

\dissertationSubheading[sec:map-h2]{H2 / RQ2: Same-window alignment}

\begin{itemize}
\item Statement: in the same window, each score agrees with the checks built from the edges it walks on.
\item Evidence: the family means of every score with their CIs (Table~\ref{tab:alignment-hybrid}); the same-window coefficients with the wallets outside a graph kept as isolated nodes (Table~\ref{tab:tie-sensitivity}).
\item Note: the agreement is partly built in, since the proxies come from the same events as the graphs. EndorseRank's allowance and transfer coefficients have different ceilings (Section~\ref{sec:same-window-alignment}), so the claim says nothing about which of the two it follows more closely, and the isolated-node rule was tried after the results were known. None of these comparisons was fixed in advance. The paired differences between EndorseRank and AWP and their per-proxy results are in Appendix~\ref{ch:appendix-er-awp}.
\end{itemize}

\dissertationSubheading[sec:map-h3]{H3 / RQ3: Computational cost}

\begin{itemize}
\item Statement: at the same $n$, a score's computation time follows the number of edges in its graph, so EndorseRank, on the sparsest graph, runs fastest.
\item Evidence: Table~\ref{tab:benchmark-runtime} (runtime, SD, memory, iterations, $|V|$, $|E|$); Tables~\ref{tab:benchmark-scaling-incohort} and~\ref{tab:benchmark-scaling}; Figure~\ref{fig:runtime-scaling}.
\item Note: the expanded-pool stages add sampled addresses but no edges, and hence no solver-graph nodes, beyond the matched-cohort graph (Section~\ref{sub:runtime-scaling}); the claim is confined to edge count.
\end{itemize}

\dissertationSubheading[sec:map-rq4]{RQ4: Temporal holdout}

\begin{itemize}
\item Statement: with scores fixed at $t_1$, the scores predict the new owner--spender approvals and new transfer senders that appear after $t_1$ and up to $t_2$ on the spender cohort ($n=1{,}335$), with the edge type deciding which label a score ranks well, but no PageRank predicts either label better than the raw $t_1$ degree of its own edge type.
\item Evidence: Tables~\ref{tab:holdout-spenders} and~\ref{tab:holdout-diff} (Section~\ref{sec:holdout-results}); the June--August window, Table~\ref{tab:fresh-posthoc}; the spenders that had received a transfer before the freeze, Table~\ref{tab:holdout-receiving}.
\item Qualification: both labels are the constructs as observed in a later period (new approvals, new senders), not credit or trading outcomes, and both are counts of arrivals. RQ4 replaced the proposal's trading-success validation (Section~\ref{sub:gf-pr}). Apart from the increment of C-PR ($\lambda=1$), none of these comparisons was fixed in advance. The differences between EndorseRank and AWP are in Appendix~\ref{ch:appendix-er-awp}.
\end{itemize}

\dissertationSubheading[sec:map-rq5]{RQ5: Coupled operator}

\begin{itemize}
\item Statement: under the rule of 14~September, a per-node convex mixture of the raw-amount allowance and transfer layers (C-PR) does not beat both layers walked alone, and the seeded ablation (S-PR) reproduces the transfer walk. The registered replication does not confirm that C-PR improves on the transfer walk alone, as it did in an exploratory comparison on the spring holdout: F1 and F3 hold, F2 fails. With AWP as the comparator, F2 fails by a wider margin.
\item Evidence, rule of 14~September: Table~\ref{tab:holdout-diff} for the primary part, and Tables~\ref{tab:alignment-hybrid} and~\ref{tab:tau-diff-hybrid} for the secondary part (Section~\ref{sec:coupled-results}).
\item Evidence, transfer walk alone: the exploratory spring contrast in Table~\ref{tab:holdout-diff}, and the registered decision (F1--F3, with F4--F6 reported beside it) in Table~\ref{tab:fresh-contrasts} (Section~\ref{sub:fresh-results}).
\item Evidence, AWP: the registered sensitivity analysis in Table~\ref{tab:fresh-contrasts}, and the spring contrasts computed after the labels were known in Table~\ref{tab:holdout-diff}.
\item Evidence, cost: runtime rows of Table~\ref{tab:benchmark-runtime}.
\item Qualification: tested with fixed $\lambda\in\{0.25,0.5,0.75\}$, raw-amount layers, uniform teleportation, transition-level mixing and one chain. The rule of 14~September and the values of $\lambda$ were set in the configuration before any coupled score was computed (Section~\ref{sec:holdout-protocol}); the rule for the comparison with the transfer walk alone was written on 28~September~2026 for the June--August window and registered on 1~October~2026 together with two sensitivity analyses, one with AWP as the comparator, before the extraction (Section~\ref{sub:fresh-holdout}). The registration documents call the two layers walked alone EndorseRank and AWP.
\end{itemize}

\dissertationSubheading[sec:map-c1c4]{Claims C1--C5}

\begin{itemize}
\item C1: the runtime evidence listed under H3/RQ3; the profiling run and the edge ratio are in Section~\ref{sec:benchmark-results} and Appendix~\ref{sec:pr-sparsity}. The claim is specific to EndorseRank, since C-PR costs at least as much as AWP.
\item C2: the same-window agreement of every score with the checks of its own edges (Table~\ref{tab:alignment-hybrid}), partly built in and, for EndorseRank, set by the rule for wallets outside the graph (Table~\ref{tab:tie-sensitivity}).
\item C3: the rank divergence listed under H1/RQ1 (Table~\ref{tab:tie-sensitivity}, Figure~\ref{fig:rank-scatter}).
\item C4: the comparison of every PageRank with the raw degree of its own edge type on the spender holdout (Tables~\ref{tab:holdout-spenders} and~\ref{tab:holdout-diff}), with the coverage check (Table~\ref{tab:holdout-receiving}), the exploratory trader run as its limit on the traders (Table~\ref{tab:holdout-traders}) and the June--August window as its repetition (Table~\ref{tab:fresh-posthoc}); the robustness checks and the sample-definition sensitivity are reported in Tables~\ref{tab:robustness-damping}--\ref{tab:robustness-sample-size}, Table~\ref{tab:endorserank-restarts} and Figure~\ref{fig:damping-sensitivity}.
\item C5: coupling does not beat its two layers walked alone, a bounded negative result, and the registered replication does not confirm its exploratory gain over the transfer walk alone (Table~\ref{tab:fresh-contrasts}); see RQ5 above and Section~\ref{sec:coupled-results}.
\end{itemize}

\dissertationSubheading[sec:map-nonclaims]{Explicit non-claims}

This research does \emph{not} claim:

\begin{itemize}
\item that any of the scores is the better measure of reputation in general, or that AWP has become obsolete (the comparisons between EndorseRank and AWP in Appendix~\ref{ch:appendix-er-awp} follow the constructs and depend on coverage);
\item that EndorseRank aligns more closely with the allowance family than with the transfer family (the two coefficients have different ceilings);
\item that any PageRank of this thesis forecasts the future growth of its own construct better than the raw $t_1$ degree does;
\item that allowance edges lead transfer edges on liquidation avoidance at the PageRank level, or that EndorseRank beats AWP there (the PageRank contrasts of Section~\ref{sub:neutral-label-results} are not resolved);
\item that any score measures creditworthiness, default risk or trading skill (the allowance signal is lower on the profitable share, and the data contain no default label);
\item that the liquidation-avoidance signal belongs to the allowance edge and not to account type (almost every approval-receiving trader is a contract account, and account type is unresolved);
\item that the value parameter $b=1$, or EndorseRank's uniform restart, is the best choice (both were fixed by this study, the restart after comparing the two rules on these data);
\item that C-PR improves on a walk over the transfers alone (the registered rule failed on new senders) or on AWP (it falls behind AWP on new senders in both windows);
\item that no combination of the two edge types can beat the single-layer methods, only that the per-node convex mixture tested here, with fixed $\lambda$, does not;
\item that runtime scales favorably when the edge population grows beyond the cohort graph;
\item that, without replication, the results generalize beyond Arbitrum One and the study window;
\item that on-chain scores substitute for collateral or legal identity systems.
\end{itemize}

\dissertationSubheading[sec:map-artifacts]{Replication artifacts}

\begin{enumerate}
\item A merged wallet-ranking table with the scores of EndorseRank, AWP, C-PR ($\lambda\in\{0,0.25,0.5,0.75,1\}$) and S-PR.
\item The machine-readable holdout summaries for the spender cohort and the matched traders, with the paired contrasts and primary-criterion verdict (committed under \texttt{data/3-published-results-for-thesis/spring-holdout-2026-03-to-2026-05/}).
\item The machine-readable evaluation summary: authoritative same-window numbers, per-run timings and the bootstrap protocol (committed under \texttt{data/3-published-results-for-thesis/}).
\item The registered replication: the registration file and plan as committed on 1~October~2026, the extraction manifest with the hashes of both files, the evaluation summary with every coefficient, contrast and verdict, and the post hoc scores of EndorseRank on the same cohorts (committed under \texttt{config-copies/} and \texttt{registered-replication-2026-06-to-2026-08/} in \texttt{data/3-published-results-for-thesis/}).
\item The exported table fragments (\texttt{results/tables/}).
\item The exported vector figures (\texttt{Figures/generated/}).
\item The extracted logs and the tables built from them (\texttt{data/1-raw-blockchain-logs/} and \texttt{data/2-processed-tables-and-evaluations/}); files over 100\,MiB are left out and can be rebuilt from the monthly files.
\item The archived extended-baseline matrices (\texttt{data/4-archived-extended-baselines/}, optional).
\end{enumerate}

The tables and figures are rebuilt by running the evaluation pipeline on real data with the robustness sweeps, exporting the tables, and then redrawing the figures from the same artifacts (Appendix~\ref{sec:pipeline-stages}).

\dissertationSubheading[sec:map-c6]{C6 --- Outcomes built from neither edge type}

\begin{itemize}
\item Statement: on the trader outcome built from neither edge type that bears on counterparty risk, later liquidation avoidance, the approval count at the freeze ranked traders above the transfer count in the registered window, after stratification by the number of closes, and in the spring window.
\item Evidence: Tables~\ref{tab:neutral-label-contrasts} and~\ref{tab:neutral-label-recipients} (Section~\ref{sub:neutral-label-results}); all six trader labels in both windows, Table~\ref{tab:neutral-label-grid}; the protocol, Section~\ref{sub:neutral-label-protocol}; the plan and claim ledger, \texttt{analysis/neutral-label/PLAN.md}, committed in \texttt{20e250f} before the analysis ran, with its script and results in the same folder.
\item Qualification: a pre-specified post hoc analysis of existing data that adds no research question. The claim is count-level, since the like-for-like PageRank pair points the same way but is not resolved. Most recipients are contract accounts, so account type is an unresolved confound, and the same signal is lower on the profitable share, so it ranks forced-liquidation avoidance, not trading success. Any use in screening would first need a registered test on new data with account type controlled (Section~\ref{sec:future-work}).
\end{itemize}

% !TEX root = ../main.tex
\section{Ethics Clearance Certificate}
\label{ch:appendix-ethics}

\dissertationSubheading[sec:ethics-certificate]{Certificate}

This study uses public blockchain data only and involves no human participants (Section~\ref{sec:ethics}). The ethics clearance certificate issued by the Unisa School of Computing / College of Agriculture and Environmental Sciences research ethics review committee is reproduced below.

% TODO (user action): replace the placeholder box with the certificate PDF once
% the reference number, date and scanned document are available, e.g.
%   \includegraphics[width=\textwidth,page=1]{Figures/ethics-clearance.pdf}
% Do not invent a reference number or date.
\begin{center}
\fbox{\parbox{0.9\textwidth}{\centering\vspace{2em}
Ethics clearance certificate\\[0.5em]
Reference: [Unisa SoC/CAES ethics clearance ref: TBD]\\
Date of issue: [TBD]\\[0.5em]
(Scanned certificate to be inserted here.)
\vspace{2em}}}
\end{center}

